\documentclass[11pt]{article}
\usepackage[margin=1in]{geometry}
\usepackage{amsmath,amssymb,amsthm}
\usepackage{mathtools}
\usepackage{graphicx}
\usepackage{booktabs}
\usepackage{multirow}
\usepackage{natbib}
\usepackage{hyperref}
\usepackage{enumitem}
\usepackage{xcolor}
\usepackage[affil-it]{authblk}
\usepackage{setspace}
\usepackage{caption}
\onehalfspacing

\captionsetup[table]{skip=8pt}
\hypersetup{colorlinks=true, linkcolor=blue!50!black, citecolor=blue!50!black, urlcolor=blue!60!black}

\graphicspath{{../figures/}}

\theoremstyle{plain}
\newtheorem{theorem}{Theorem}
\newtheorem{lemma}[theorem]{Lemma}
\newtheorem{corollary}[theorem]{Corollary}
\newtheorem{proposition}[theorem]{Proposition}

\theoremstyle{definition}
\newtheorem{definition}[theorem]{Definition}
\newtheorem{example}[theorem]{Example}
\newtheorem{assumption}[theorem]{Assumption}

\theoremstyle{remark}
\newtheorem{remark}[theorem]{Remark}
\newtheorem{observation}[theorem]{Observation}

\newcommand{\F}{\mathcal{F}}
\newcommand{\N}{N}
\newcommand{\CW}{\mathrm{CW}}
\newcommand{\cstar}{c^{\ast}}

\title{\textbf{Online Traveling Salesman Games: \\ Temporal Nucleolus and the Fragility of the Core under Dynamic Arrivals}}

\author[1]{Seyun Jeong}
\author[2]{Hyunchul Tae}
\affil[1]{Department of Industrial and Management Engineering, Pohang University of Science and Technology (POSTECH), 77 Cheongam-ro, Nam-gu, Pohang, Gyeongbuk 37673, Republic of Korea}
\affil[2]{Korea Institute of Industrial Technology (KITECH), 89 Yangdaegiro-gil, Ipjang-myeon, Seobuk-gu, Cheonan-si, Chungcheongnam-do 31056, Republic of Korea}

\date{\today}

\begin{document}
\maketitle

\section*{Highlights}
\begin{itemize}
\item Online TSG: coalitions restricted to feasibility under arrival dynamics.
\item Temporal Nucleolus via sequential LP over the feasibility family $\F$.
\item Sufficient emptiness bound $r>r^{\ast\ast}$, tight for complement coalitions.
\item Structural safeguard $k<n-1$ blocks the complement-coalition mechanism.
\item Asymptotic tightening $r^{\ast\ast}\to 1$ at $O(n^{-1/2})$ confirmed to $n=50$.
\end{itemize}

\begin{abstract}
\noindent We introduce the \emph{Online Traveling Salesman Game} (Online TSG), a cooperative game in which customers arrive over time and cost must be allocated only among coalitions realizable under the revealed arrival dynamics. The classical static TSG is recovered at simultaneous arrivals. We formalize the family $\F$ of feasible coalitions, define the Temporal Nucleolus as the lexicographic minimizer of excess restricted to $\F$, and give a sequential LP procedure to compute it. Our theoretical results address Core stability under dynamics through two complementary sufficient conditions for emptiness: (i) $r>r^{\ast\ast}$, with $r^{\ast\ast}=1+\min_i \delta_i/\cstar(N)$ (Theorem~\ref{thm:empty-core}, single-complement mechanism); and (ii) $r>r^{\ast\ast\ast}$, with $r^{\ast\ast\ast}=\frac{1}{n-1}\sum_i c(N\setminus\{i\})/\cstar(N)$ (Proposition~\ref{prop:bondareva-complement}, Bondareva--Shapley balanced-complement threshold). A bounded peak queue $k<n-1$ structurally blocks both complement-based mechanisms (Corollary~\ref{cor:cor52}); and $r^{\ast\ast}\to 1$ almost surely at rate $O(n^{-1/2})$ (Theorem~\ref{thm:asymptotic}), tying the cooperative-game threshold to the Beardwood--Halton--Hammersley $\sqrt{n}$ scaling of the Euclidean TSP. Over 175 random instances with $n\in\{5,7,10,12,15\}$ and seven arrival patterns, empty Cores decompose into four disjoint mechanisms: Theorem~\ref{thm:empty-core} fires in $37$ instances, Proposition~\ref{prop:bondareva-complement} adds $11$ balanced-complement cases, a residual $10$ near-complement cases (Remark~\ref{rem:near-complement}) and $9$ intermediate-coalition cases (Observation~\ref{obs:intermediate}) complete the $67$ observed empties; no false positives arise under any of three dispatch policies. The complement-mechanism thresholds $\bar r^{\ast\ast}=1.223$ and $\bar r^{\ast\ast\ast}=1.096$ are sharper than the classical bound $\bar r^{\ast}=4.351$ by factors of $3.56$ and $3.97$ respectively. Scale-up experiments to $n=50$ confirm the asymptotic decay, and the intermediate mechanism is observed exclusively under nearest-neighbor dispatch, suggesting that short service windows together with moderately-informed dispatch policies are sufficient to preserve Core stability.
\end{abstract}

\noindent\textbf{Keywords:} Game theory; Online optimization; Cooperative game theory; Traveling salesman game; Nucleolus; Core stability; Dynamic routing.

\section{Introduction}
\label{sec:introduction}

Urban last-mile logistics is characterized by requests that enter the system \emph{after} planning begins: online grocery orders, same-day parcels, ride-hail queues and crowdsourced pickups all share the feature that the identity of who will be served together with whom is not known until dispatch time \citep{pillac2013review,ausiello2001oltsp}. Operators nevertheless must share the routing cost among the customers who are eventually served on a joint tour, and the quality of this allocation shapes incentives for participation and pricing \citep{guajardo2016review}. Cooperative game theory, and specifically the Traveling Salesman Game (TSG) of \citet{potters1992tsg}, supplies the canonical framework in which a tour cost is divided so that no subset of customers has a strict incentive to defect; the Vehicle Routing Game (VRG) of \citet{gothe1996vrg,engevall2004heterogeneous} extends the idea to multi-vehicle settings. Classic applications of the VRG include school bus routing \citep{park2012schoolbus}, where the set of customers is fixed before the vehicle departs. Online counterparts arise naturally in same-day delivery and ride-hailing platforms, where requests arrive dynamically and the coalition of customers to be served jointly is revealed only during execution.

The classical TSG, however, is defined under the tacit assumption that every customer is available to be visited at any point of the tour: any coalition $S\subseteq N$ can be \emph{imagined} as a self-contained instance, and its characteristic cost $c(S)$ is the optimal TSP tour on $S\cup\{0\}$. In online settings, this counterfactual is not available: a coalition $S$ is operationally meaningful only if there exists a moment at which all members of $S$ have arrived but none have yet been served. Ignoring temporality hides a central source of instability. Two customers who arrive far apart in time are never candidates to be served together, and pretending that they form a coalition distorts fairness. Recent work on online cooperative games \citep{ge2024incentives,aziz2025participation,goyal2025temporal} has begun to address temporality but not in the routing context.

\paragraph{Research question.} What is the correct cooperative game-theoretic abstraction of the online traveling salesman problem, how does it compute, and---most importantly---when does it admit a stable allocation?

\paragraph{Contributions.} We make four contributions.
\begin{itemize}[leftmargin=1.5em]
\item[\textbf{C1.}] \emph{Model.} We define the Online TSG $(N, c, \F)$ by restricting coalitional stability constraints to the family $\F$ of proper sub-coalitions realizable under the realized arrival dynamics, with the grand coalition handled through the efficiency equality. This places the Online TSG within the taxonomy of restricted-cooperation games (Section~\ref{sec:related}) but with a novel source of restriction: $\F$ is not an exogenous graph or combinatorial structure but an endogenous random family induced by the arrival-service process. We introduce the Temporal Nucleolus as the lexicographic minimizer of excess over this time-induced restricted family. The classical static TSG is recovered as the special case where arrivals collapse to a single instant and the realized tour is optimal (Theorem~\ref{thm:static-equiv}).
\item[\textbf{C2.}] \emph{Core theory under dynamics.} We provide two complementary sufficient conditions for empty Core: Theorem~\ref{thm:empty-core} ($r>r^{\ast\ast}$, single-complement) and Proposition~\ref{prop:bondareva-complement} ($r>r^{\ast\ast\ast}$, balanced-complement via Bondareva--Shapley on the collection of all complements). Corollary~\ref{cor:cor52} provides a structural obstruction: a bounded peak queue $k=\max_t|U_t|<n-1$ blocks both complement-based mechanisms. Together these delineate the regimes in which complement-coalition instability arises. Two residual mechanisms remain empirically active: a near-complement mechanism (Remark~\ref{rem:near-complement}) mixing $(n-1)$- and $(n-2)$-coalitions, and an intermediate-coalition mechanism (Observation~\ref{obs:intermediate}) operating at $k<n-1$ with binding on intermediate-size coalitions; analytic characterizations of both are left open.
\item[\textbf{C3.}] \emph{Asymptotic refinement.} Under uniform coordinate sampling, $r^{\ast\ast}\to 1$ at rate $O(n^{-1/2})$ almost surely (Theorem~\ref{thm:asymptotic}); this explains the empirical smallness of $r^{\ast\ast}$ without reference to specific simulation parameters and ties the cooperative-game threshold to the classical Beardwood--Halton--Hammersley $\sqrt{n}$ scaling of the Euclidean TSP.
\item[\textbf{C4.}] \emph{Empirical calibration.} On 175 simulated instances across five coalition sizes and seven arrival patterns under three dispatch policies (525 policy-instance pairs), plus 45 scale-up instances up to $n=50$ and a 24-instance scale-invariance study (Appendix~\ref{app:scale-invariance}), we verify every theoretical claim. Under NN dispatch, Theorem~\ref{thm:empty-core} fires in 37 instances (all observed empty) and Proposition~\ref{prop:bondareva-complement} adds 11 balanced-complement certifications; a further 10 near-complement and 9 intermediate-coalition empty Cores complete the 67 observed empties. No false positives arise under any policy. The mean thresholds ($\bar r^{\ast\ast}=1.223$, $\bar r^{\ast\ast\ast}=1.096$) are sharper than the classical bound ($\bar r^{\ast}=4.351$) by factors of $3.56$ and $3.97$. Corollary~\ref{cor:cor52} holds analytically without exception: no instance with $k<n-1$ satisfies either threshold; the 9 intermediate-coalition empties all lie at $k<n-1$ under NN and vanish under cheapest-insertion and batch-reoptimization. The asymptotic decay $r^{\ast\ast}\to 1$ is consistent with Theorem~\ref{thm:asymptotic}'s $O(n^{-1/2})$ rate on the scale-up grid.
\end{itemize}

\paragraph{Practical implication.} In operational terms, our results say that a bounded peak queue $k<n-1$---achievable via short service-time promises---eliminates all complement-coalition empty-Core mechanisms (Theorem~\ref{thm:empty-core} and Proposition~\ref{prop:bondareva-complement}). This covers the pathological ``morning batch'' regime where many orders sit simultaneously in the queue. In the shallow-queue regime, however, the intermediate-coalition mechanism of Observation~\ref{obs:intermediate} remains active and can empty the Core when the realized online competitive ratio is large. This mechanism is policy-sensitive: it was observed only under nearest-neighbor dispatch in our experiments and vanishes under cheapest-insertion or batch re-optimization (Section~\ref{ssec:policy}). Jointly, our findings suggest that short service windows together with moderately-informed dispatch policies are sufficient to preserve Core stability.

\paragraph{Structure.} Section~\ref{sec:related} situates our work within the literature on static TSG/VRG, dynamic and rolling-horizon games, online cooperative games, and nucleolus computation. Section~\ref{sec:problem} formalizes the Online TSG. Section~\ref{sec:nucleolus} defines and computes the Temporal Nucleolus. Section~\ref{sec:theory} develops the Core theory. Section~\ref{sec:experiments} reports computational experiments, and Section~\ref{sec:conclusion} concludes.

\section{Related Work}
\label{sec:related}

We organize prior work along four streams and state, for each, the precise point of departure of our contribution.

\subsection{Static Traveling Salesman and Vehicle Routing Games}
The cooperative TSP was initiated by \citet{potters1992tsg} with a characteristic function equal to the optimal TSP tour cost on the chosen coalition. \citet{tamir1989core} showed that the Core may be empty even on the Euclidean line, and the gap between fair and feasible allocations was quantified via approximation in \citet{faigle1998approximate} and \citet{blaser2008metric}. In parallel the Vehicle Routing Game (VRG) received a Nucleolus-based treatment by \citet{gothe1996vrg}, extended to heterogeneous fleets by \citet{engevall2004heterogeneous} and to time-window settings by \citet{tae2020vrgtw} and \citet{arroyo2024vrptw}. The Shapley value formulation of cooperative TSP under rolling horizons is treated by \citet{kimms2016rolling}, and a core-based variant by \citet{kimms2016core}. A broad survey is given by \citet{guajardo2016review}. \emph{Departure.} Every cited study defines coalitions over the full customer set $2^N$: coalition identity is assumed known \emph{a priori}. We replace $2^N$ with the data-dependent family $\F$ of coalitions realizable under the online arrival process.

\subsection{Dynamic and Rolling-Horizon Cooperative Games}
\citet{kranich2005dynamic} proposed Core concepts for TU games played over discrete time. \citet{habis2010weak} defined the weak sequential Core to accommodate renegotiation, and \citet{lehrer2013dynamic} studied how the Core evolves as information unfolds. \citet{kimms2016rolling} applied rolling horizons to cooperative TSP, redistributing cost as new customers enter. \emph{Departure.} These works abstract away the physical routing cost structure: coalitions are described combinatorially with exogenous values $v(S,t)$. In contrast we retain the metric TSP cost $c(S)$ and let temporality enter through the \emph{admissibility} of coalitions, which couples the game to the dispatch policy.

\subsection{Online Cooperative Games}
\citet{bullinger2023online} studied online coalition formation under random arrival and dissolution. \citet{ge2024incentives} and \citet{zhang2025incentives} designed incentives for early arrival in online cost-sharing settings. \citet{aziz2025participation} characterized participation incentives for players entering over time. \citet{goyal2025temporal} introduce Temporal Cooperative Games, in which the characteristic function is indexed by time. \emph{Departure.} These models differ from ours in two respects. First, \citet{bullinger2023online} formalize coalition \emph{formation} (which players group with whom) rather than cost sharing within a fixed grand coalition, so their stability concepts and ours address distinct strategic problems. Second, and common to the whole stream, none of these models equip the online coalition process with a TSP or routing cost structure; the inequalities that define our Temporal Core reflect exactly the geometric interactions among customers in the plane, allowing the sharp thresholds $r^{\ast\ast}$ and structural bound $k<n-1$ established in Section~\ref{sec:theory}.

\subsection{Nucleolus Computation}
\citet{schmeidler1969nucleolus} introduced the Nucleolus as the lexicographic minimizer of the excess vector. The sequential-LP approach that we adopt descends from Kohlberg's algorithm and its careful modernization by \citet{guajardo2015nucleolus}, with further scaling to large games by \citet{nguyen2016large}. Nucleolus-based ridesharing is treated in \citet{lu2019ridesharing}, and Nucleolus computation for the VRG with time windows (VRPTW) in \citet{tae2020vrgtw}. \emph{Departure.} Our LP cascade is defined over $\F$ rather than $2^N\setminus\{\emptyset\}$: in sequential arrival regimes this can reduce the per-LP constraint count by more than an order of magnitude (Section~\ref{sec:experiments}), keeping the method computationally viable for realistic problem sizes.

\subsection{Restricted Cooperation Games}

A parallel line of cooperative game theory, independent of the routing literature, studies games in which the set of formable coalitions is a \emph{strict subset} of $2^N$. \citet{myerson1977graphs} initiated this theme by treating coalitions as paths in an underlying communication graph, producing the graph-restricted game $(N, v^{c})$ in which $v^{c}(S)$ aggregates values of $S$'s connected components. \citet{faigle1989cores} generalized the framework to games on combinatorial structures (antimatroids, convex geometries, partitions), and \citet{bilbao2000} provided a unified treatment. \citet{algaba2001unified} connected Myerson's and Faigle's lines. More recent extensions handle permission structures \citep{gilles2010economic}, hierarchical coalition configurations, and the core of games on ordered structures \citep{grabisch2013core}. Computing the nucleolus of a restricted game is in general no harder than the unrestricted case and may be substantially easier, since the constraint count is $|\F|$ rather than $2^n-1$ \citep{schouten2019nucleolus}. \emph{Departure.} In this restricted-cooperation literature, feasibility is determined by an \emph{exogenous and time-invariant} structure---a graph, a partition, an antimatroid, or a permission relation---specified as part of the game's primitives. The Online TSG introduces a different source of restriction: feasibility is \emph{endogenously induced} by the realized dynamics of a stochastic arrival-service process. Formally, $\F$ is a data-dependent random family with $S\in\F$ iff customers in $S$ share a nonempty common waiting window $\CW(S)=\bigcap_{i\in S}[a_i,s_i)$ (Definition~\ref{def:feasible}). The resulting family $\F$ is closed under neither union nor intersection in general, and its shape depends on the arrival-order permutation, the queue-depth process, and the dispatch policy. Our theoretical contributions---the complement-coalition threshold $r^{\ast\ast}$ (Theorem~\ref{thm:empty-core}), the structural obstruction at bounded peak queue (Corollary~\ref{cor:cor52}), and the asymptotic tightening (Theorem~\ref{thm:asymptotic})---are specific to this time-induced restriction and have no direct counterpart in the graph-restricted or antimatroid literatures. We view the Online TSG as adding \emph{temporally-induced restricted cooperation} to the taxonomy of restricted-cooperation games.

A separate line of theoretical work relevant to our asymptotic analysis (Theorem~\ref{thm:asymptotic}) is the probabilistic theory of Euclidean combinatorial optimization, initiated by \citet{beardwood1959shortest} and systematized by \citet{steele1981subadditive,steele1997probability}. The Beardwood--Halton--Hammersley theorem characterizes the $\sqrt{n}$ scaling of the optimal Euclidean TSP tour, which in our setting controls the denominator of $r^{\ast\ast}$ and hence its asymptotic tightening. To the best of our knowledge, this asymptotic perspective on cooperative routing games is new: prior cooperative TSP and VRG studies \citep{potters1992tsg,gothe1996vrg,engevall2004heterogeneous,kimms2016core} treat the coalition cost as a fixed quantity rather than as a random variable with a known limit distribution.

\section{Problem Definition}
\label{sec:problem}

\subsection{Notation}
We consider a depot $0$ and $n$ customers $N = \{1,\ldots,n\}$, located in a metric space with symmetric travel cost $d$. A single vehicle leaves the depot, visits a subset of customers, and returns. Each customer $i\in N$ is associated with a random \emph{arrival time} $a_i \in \mathbb{R}_+$, revealed online, at which its request enters the system, and a \emph{service time} $s_i \ge a_i$ at which the vehicle first serves it in the realized execution. Table~\ref{tab:notation} collects the notation used throughout.

\begin{table}[t]
\centering
\caption{Notation.}
\label{tab:notation}
\small
\begin{tabular}{@{}ll@{}}
\toprule
Symbol & Meaning \\
\midrule
$N,\; n$ & Grand coalition of customers and its cardinality \\
$d(i,j)$ & Symmetric travel cost (metric) from $i$ to $j$ \\
$c(S)$ & Optimal TSP tour cost on $S\cup\{0\}$, \emph{time-independent} \\
$\cstar(N)$ & $c(N)$; the grand-coalition TSP optimum \\
$a_i,\; s_i$ & Arrival and service time of customer $i$ \\
$U_t$ & Customers that have arrived by $t$ and are not yet served, $U_t=\{i:a_i\le t< s_i\}$ \\
$\F$ & Family of \emph{feasible coalitions}: $S\in\F \iff \exists\, t,\ S\subseteq U_t$ \\
$\CW(S)$ & Common window $\bigcap_{i\in S}[a_i, s_i)$ \\
$C(N)_{\text{online}}$ & Cost of the online tour realized by the dispatch policy \\
$r$ & Competitive ratio $C(N)_{\text{online}}/\cstar(N)$ \\
$\delta_i$ & $c(\{i\})+c(N\setminus\{i\})-\cstar(N)$ \\
$r^{\ast\ast}$ & $1+\min_{i:\,N\setminus\{i\}\in\F}\delta_i/\cstar(N)$: tight threshold (Thm.~\ref{thm:empty-core}); undefined if no such $i$ exists \\
$r^{\ast}$ & $\sum_i c(\{i\})/\cstar(N)$: Remark~\ref{rem:rstar} upper bound \\
$k$ & $\max_t |U_t|$: peak size of the waiting queue \\
\bottomrule
\end{tabular}
\end{table}

\subsection{Feasible coalitions and the characteristic function}

The starting point is that the characteristic cost $c(S)$ of a coalition is \emph{defined in the same way as in the static TSG}: it is the cost of the optimal TSP tour on $S\cup\{0\}$ irrespective of time. Temporality enters only through the admissibility of coalitions.

\begin{definition}[Common waiting window]
For a coalition $S\subseteq N$,
$\CW(S) := \bigcap_{i\in S} [a_i, s_i)$.
\end{definition}

\begin{definition}[Feasible coalition]
\label{def:feasible}
A nonempty coalition $S\subseteq N$ is \emph{feasible} under the realized arrival/service sequence $(a, s)$ if $\CW(S)\neq\emptyset$, equivalently if there exists an event time $t$ at which $S\subseteq U_t$. Let
$\F := \{S\subseteq N : \CW(S)\neq\emptyset\}$
denote the family of feasible coalitions.
\end{definition}

Note that $\{i\}\in\F$ for every $i$ (singletons are always feasible at $a_i$), and $N\in\F$ whenever all $n$ customers sit simultaneously in the queue at some instant; however, $N$ appears in the Temporal Core only through the efficiency equality (Definition~\ref{def:online-tsg}) and not as a stability constraint.

\begin{definition}[Online TSG]
\label{def:online-tsg}
An \emph{Online Traveling Salesman Game} is a triple $(N, c, \F)$ where $N$ is the customer set, $c:2^N\setminus\{\emptyset\}\to\mathbb{R}_+$ is the TSP characteristic cost defined on every nonempty subset (and used in core constraints only for $S\in\F$), and $\F\subseteq 2^N\setminus\{\emptyset\}$ is the family from Definition~\ref{def:feasible}.
\end{definition}

An \emph{allocation} is a vector $x\in\mathbb{R}^n$ with $\sum_i x_i = C(N)_{\text{online}}$, where $C(N)_{\text{online}}$ denotes the realized online tour cost. Following the standard convention in restricted-cooperation cooperative games \citep{myerson1977graphs,bilbao2000}, the grand coalition $N$ is handled through the efficiency equality, and coalitional stability constraints are imposed only on \emph{proper} feasible sub-coalitions. The \emph{Temporal Core} is
\[
\mathrm{Core}(N,c,\F) := \bigl\{ x\in\mathbb{R}^n : \textstyle\sum_i x_i = C(N)_{\text{online}},\; \sum_{i\in S} x_i \le c(S)\ \forall S\in\F,\ S\subsetneq N \bigr\}.
\]
When $C(N)_{\text{online}} = \cstar(N)$, this definition coincides with the classical \citet{potters1992tsg} static TSG Core on the same constraint family. In general, the online overhead $C(N)_{\text{online}}-\cstar(N)\ge 0$ raises the efficiency budget above the minimum grand-coalition cost; Table~\ref{tab:static-vs-temporal} summarizes the contrast, and the competitive ratio $r$ (defined below) quantifies this overhead.

\paragraph{Characteristic cost versus realized cost.} The characteristic function $c(S)$ is the time-independent optimal TSP tour on $S\cup\{0\}$, matching the classical \citet{potters1992tsg} definition. The realized online tour cost $C(N)_{\text{online}}$, by contrast, is the total travel distance of the tour actually executed by the online dispatch policy, which may exceed $\cstar(N)$ because the online policy cannot anticipate future arrivals; in general $C(N)_{\text{online}} \ge \cstar(N)$. We use $\sum_i x_i = C(N)_{\text{online}}$ as the efficiency budget because this is the cost the platform must actually allocate; the competitive ratio $r = C(N)_{\text{online}}/\cstar(N)$ measures the overhead introduced by online dispatch, since the online policy must commit to visits without seeing future arrivals. The threshold $r^{\ast\ast}$ (Theorem~\ref{thm:empty-core}) is a function of coordinate geometry alone through $c(\{i\})$, $c(N\setminus\{i\})$, and $\cstar(N)$.

\begin{table}[t]
\centering
\caption{Static TSG versus Online TSG.}
\label{tab:static-vs-temporal}
\small
\begin{tabular}{@{}lll@{}}
\toprule
 & Static TSG \citep{potters1992tsg} & Online TSG (this paper) \\
\midrule
Coalitions considered & All $S\subseteq N$, $S\neq\emptyset$ & $S\in\F$ \\
Characteristic function & $c(S)$ = TSP on $S\cup\{0\}$ & Same \\
Core constraints & $2^n-1$ & $|\F|\le 2^n-1$ \\
Recovers classical TSG & --- & Yes, when $\F=2^N\setminus\{\emptyset\}$ \\
\bottomrule
\end{tabular}
\end{table}

\begin{example}[Solomon C101, first five customers]
\label{ex:solomon}
We illustrate the Online TSG on the first five customers of the Solomon C101 benchmark \citep{solomon1987benchmarks}. Depot is at $(35, 35)$ and coordinates/ready times are as in Table~\ref{tab:solomon}. The online policy is nearest-neighbor departing from the depot at $t=0$.

\begin{table}[h]
\centering
\caption{Solomon C101 first five customers used in Example~\ref{ex:solomon}.}
\label{tab:solomon}
\small
\begin{tabular}{@{}cccccc@{}}
\toprule
$i$ & $(x_i, y_i)$ & ready $a_i$ & due & $c(\{i\})$ & service $s_i$ \\
\midrule
1 & (45, 68) & 912 & 967 & 68.964 & 912.00 \\
2 & (45, 70) & 825 & 870 & 72.801 & 825.00 \\
3 & (42, 66) &  65 & 146 & 63.561 &  65.00 \\
4 & (42, 68) & 727 & 782 & 67.469 & 727.00 \\
5 & (40, 69) &  15 &  67 & 68.731 &  34.37 \\
\bottomrule
\end{tabular}
\end{table}

Euclidean distances give $\cstar(N)=77.593$ and realized online cost $C(N)_{\text{online}}=80.059$, so $r = 1.032$. Because the vehicle spends most of its time idle waiting for each successive ready window, the service timeline is fully sequential in the sense of Lemma~\ref{lem:lemma52}: one verifies that $s_j \le a_{j+1}$ for every consecutive pair in the ready-time order $(5,3,4,2,1)$. Consequently, by Corollary~\ref{cor:cor51} (equivalently Corollary~\ref{cor:cor52}, since $k = \max_t|U_t| = 1 < n-1=4$), Theorem~\ref{thm:empty-core} is vacuous: \emph{no} complement coalition $N\setminus\{i\}$ is feasible. The quantities $\delta_i = c(\{i\}) + c(N\setminus\{i\}) - \cstar(N)$ are still well-defined as numerical combinations of TSP costs---one obtains $\delta_i \in [62.66, 70.20]$ and hence the formal ratio $1+\min_i \delta_i/\cstar(N) \approx 1.81$ (this numerical value is shown for illustrative comparison only; $r^{\ast\ast}$ itself is undefined in this instance because the minimum in its definition is taken over the empty set). The inequality $r > r^{\ast\ast}$ is thus not testable and Theorem~\ref{thm:empty-core} makes no prediction here.

Adopting the convention $\{i\}\in\F$ (singletons are always feasible at arrival) and using $\sum_i x_i = C(N)_{\text{online}}$ as the efficiency budget, the Temporal Nucleolus minimizes $\max_i\{x_i - c(\{i\})\}$ subject to efficiency, yielding the equal-excess allocation
\[
(x_1, x_2, x_3, x_4, x_5) \;=\; (16.67,\; 20.51,\; 11.27,\; 15.18,\; 16.44),
\]
with uniform slack $c(\{i\})-x_i = 52.29$ across customers. The static TSG allocation (computed on $2^N\setminus\{\emptyset\}$, treating every coalition as feasible) is instead the Nucleolus of the full 31-constraint game; this example shows how temporal restriction discards $26$ of those $31$ constraints while preserving core existence.
\end{example}

\section{The Temporal Nucleolus}
\label{sec:nucleolus}

Even when the Core is nonempty it is typically a polytope, and one needs a selection rule. We adopt the nucleolus of \citet{schmeidler1969nucleolus}, adapted to the feasible family $\F$.

\subsection{Definition}
For an allocation $x$ and $S\in\F$, the \emph{excess} is $e(S, x) := \sum_{i\in S} x_i - c(S)$. Let $\theta(x)$ be the vector $(e(S,x))_{S\in\F}$ sorted in non-increasing order.

\begin{definition}[Temporal Nucleolus]
\label{def:tnu}
The Temporal Nucleolus of $(N, c, \F)$ is
\[
\mathrm{TNu}(N,c,\F) := \arg\!\min{}^{\text{lex}} \bigl\{ \theta(x) : \textstyle\sum_i x_i = C(N)_{\text{online}}\bigr\},
\]
where the minimization is lexicographic over the sorted excess vector.
\end{definition}

\subsection{Sequential LP computation}

Following the standard technique of \citet{guajardo2015nucleolus}, the Temporal Nucleolus is computed by solving at most $|\F|$ linear programs. At stage $q$ the LP
\begin{equation}
\label{eq:lp-q}
\begin{aligned}
\min_{x,\varepsilon} \;& \varepsilon \\
\text{s.t. } \;& \sum_{i} x_i = C(N)_{\text{online}}, \\
& \sum_{i\in S} x_i - c(S) \le \varepsilon \qquad \forall S\in\F\setminus\F_{<q}, \\
& \sum_{i\in S} x_i - c(S) = \varepsilon_{q(S)}^{\ast} \qquad \forall S\in\F_{<q},
\end{aligned}
\end{equation}
determines the $q$-th largest coordinate of $\theta$. Here $\F_{<q}$ collects the coalitions whose excess was fixed at earlier stages with optimal value $\varepsilon_{q(S)}^{\ast}$. Finiteness of $\F$ guarantees termination in at most $|\F|$ stages; each LP has $n+1$ variables and $|\F|+1$ constraints.

\paragraph{Algorithmic details.} Correct Nucleolus computation via the sequential LP requires two subtleties beyond the basic cascade described above. First, when the optimal $\varepsilon_q^{\ast}$ is achieved simultaneously by multiple constraints (the degenerate case), simply fixing \emph{all} tight constraints at $\varepsilon_q^{\ast}$ and proceeding can fail to identify the true Nucleolus; the standard remedy, due to \citet{kopelowitz1967computation} and clarified in the practical review of \citet{guajardo2015nucleolus}, is to fix only those constraints whose left-hand sides are linearly independent relative to the currently fixed set, so that the feasible region of the next LP shrinks by exactly one dimension per iteration. Second, the procedure terminates when the feasible region becomes a singleton, which happens in at most $n-1$ iterations provided $\F$ contains a collection of proper coalitions spanning $\mathbb{R}^n/\{\text{efficiency}\}$. Our PuLP/CBC implementation identifies tight constraints via the primal optimal and checks linear independence via Gaussian elimination on the corresponding constraint coefficient vectors; tight constraints that are linearly dependent on the currently fixed set are deferred to subsequent LPs, matching the Kopelowitz--Guajardo--J\"ornsten procedure.

The restriction from $2^N\setminus\{\emptyset\}$ to $\F$ materially reduces problem size whenever arrivals are spread in time. Section~\ref{sec:experiments} reports empirically $|\F|/(2^n-1)$ as low as $0.061$ for heavily sequential patterns.

\subsection{Reduction to the static game at simultaneous arrivals}

\begin{theorem}[Reduction at simultaneous arrivals]
\label{thm:static-equiv}
Suppose $a_1 = a_2 = \cdots = a_n$ and $s_i > a_1$ for every $i\in N$, so all customers share a common waiting window. Then $\F = 2^N\setminus\{\emptyset\}$, and the Temporal Nucleolus of $(N,c,\F)$ equals the lexicographic minimizer of excess on the static TSG constraint family over proper coalitions $S\subsetneq N$ with efficiency budget $C(N)_{\text{online}}$. When the realized tour is additionally optimal (i.e.\ $C(N)_{\text{online}}=\cstar(N)$), this coincides with the classical static Nucleolus of \citet{potters1992tsg}.
\end{theorem}

\begin{proof}
By assumption there exists $t^{\star}\in[a_1,\min_i s_i)$ with $U_{t^{\star}} = N$. Every $S\subseteq N$ satisfies $S\subseteq U_{t^{\star}}$, so $\CW(S)\supseteq [a_1,\min_i s_i)\neq\emptyset$ and $\F = 2^N\setminus\{\emptyset\}$. The Temporal Core (Definition~\ref{def:online-tsg}) imposes constraints for every proper subset $S\subsetneq N$---i.e.\ the full static TSG inequality family---together with the efficiency equality $\sum_i x_i = C(N)_{\text{online}}$. The sequential-LP cascade~\eqref{eq:lp-q} thus coincides with Kohlberg's procedure for the static TSG on this constraint family, and the lexicographic minimizer is unique. When $C(N)_{\text{online}}=\cstar(N)=c(N)$, the efficiency budget matches the classical static-Nucleolus LPs and the allocation recovered is that of \citet{potters1992tsg}. In general, the Temporal Nucleolus under simultaneous arrivals is the lexicographic minimizer on the same constraint family but with a shifted efficiency budget.
\end{proof}

In short, the Online TSG is a proper generalization: the static theory is the $k=n$ extremal case where the peak queue contains everybody at once.

\section{Theoretical Results}
\label{sec:theory}

We now develop the Core theory. Section~\ref{ssec:empty} establishes conditions under which the Temporal Core is \emph{empty}; Section~\ref{ssec:survive} gives conditions under which it is guaranteed \emph{nonempty}. Section~\ref{ssec:discussion} discusses the relationships between the parameters $r$, $r^{\ast\ast}$, $k$, and the queue dynamics.

\subsection{Conditions under which the Core vanishes}
\label{ssec:empty}

\begin{remark}[Necessary super-additivity failure]
\label{rem:remark1}
If the online tour obeys $C(N)_{\text{online}} > \sum_{i\in N} c(\{i\})$, then any allocation summing to $C(N)_{\text{online}}$ violates the singleton constraints $\sum_{\{i\}} x_i \le c(\{i\})$, so the Temporal Core is empty. This is only a \emph{loose} condition; in our experiments it almost never fires because $C(N)_{\text{online}}/\cstar(N) < \sum_i c(\{i\})/\cstar(N) = r^{\ast}$ by a large margin.
\end{remark}

\begin{remark}[Remark~\ref{rem:remark1} reinterpreted as a ratio]
\label{rem:rstar}
Dividing both sides by $\cstar(N)$, Remark~\ref{rem:remark1} is equivalent to $r > r^{\ast}$, where $r^{\ast} := \sum_i c(\{i\})/\cstar(N)$. Because $r^{\ast}$ is typically in the range $3$--$4$ while the competitive ratio of reasonable online policies is close to $1$, this test is rarely tight.
\end{remark}

We sharpen Remark~\ref{rem:remark1} substantially by exploiting complement coalitions of the form $N\setminus\{i\}$.

\begin{lemma}[Feasibility of $N\setminus\{i\}$]
\label{lem:lemma51}
For each $i\in N$, $N\setminus\{i\}\in\F$ if and only if $\max_{j\neq i} a_j < \min_{j\neq i} s_j$.
\end{lemma}

\begin{proof}
($\Rightarrow$) If $N\setminus\{i\}\in\F$ there is $t$ with $N\setminus\{i\}\subseteq U_t$. Then $a_j\le t < s_j$ for every $j\neq i$, so $\max_j a_j \le t < \min_j s_j$ (both over $j\neq i$). ($\Leftarrow$) Conversely, if $\max_{j\neq i} a_j < \min_{j\neq i} s_j$, any $t$ in the open interval between them satisfies $a_j\le t<s_j$ for all $j\neq i$, so $N\setminus\{i\}\subseteq U_t$ and $N\setminus\{i\}\in\F$.
\end{proof}

\begin{lemma}[Fully sequential arrivals preclude $N\setminus\{i\}\in\F$]
\label{lem:lemma52}
If arrivals are fully sequential in the sense that $s_j \le a_{j+1}$ for all $j=1,\ldots,n-1$ (each customer is served before the next arrives), then $N\setminus\{i\}\notin\F$ for every $i\in N$.
\end{lemma}

\begin{proof}
Fix $i$ and assume customers are indexed in arrival order $1,\ldots,n$ so that $a_1 \le \cdots \le a_n$. Let $\ell^{\star}=\max\{j:j\neq i\}$ and $j^{\star}=\min\{j:j\neq i\}$ denote, respectively, the latest and earliest arrival-order indices other than $i$. By the sequential-arrival hypothesis $s_j \le a_{j+1}$ applied along the chain $j^{\star}, j^{\star}+1, \ldots, \ell^{\star}-1$, we obtain $s_{j^{\star}} \le a_{j^{\star}+1} \le s_{j^{\star}+1} \le \cdots \le a_{\ell^{\star}}$, hence $a_{\ell^{\star}} \ge s_{j^{\star}}$. Therefore $\max_{j\neq i} a_j \ge \min_{j\neq i} s_j$, violating the if-and-only-if criterion of Lemma~\ref{lem:lemma51}.
\end{proof}

\begin{theorem}[Sufficient condition for empty Temporal Core via complement coalitions]
\label{thm:empty-core}
Assume that $N\setminus\{i\}\in\F$ for some $i\in N$. Let
\[
\delta_i := c(\{i\}) + c(N\setminus\{i\}) - \cstar(N),\qquad r^{\ast\ast}:= 1 + \min_{i:\,N\setminus\{i\}\in\F} \delta_i/\cstar(N).
\]
If $r = C(N)_{\text{online}}/\cstar(N) > r^{\ast\ast}$, then the Temporal Core is empty.
\end{theorem}

\begin{proof}
Let $i^{\star}\in\arg\min_{i:\,N\setminus\{i\}\in\F} \delta_i$. Both $\{i^{\star}\}\in\F$ (singletons are always feasible) and $N\setminus\{i^{\star}\}\in\F$. Suppose, for contradiction, that $x$ lies in the Temporal Core. The two core inequalities read
\begin{align}
x_{i^{\star}} &\le c(\{i^{\star}\}), \label{eq:sing}\\
\sum_{j\neq i^{\star}} x_j &\le c(N\setminus\{i^{\star}\}).\label{eq:comp}
\end{align}
Adding~\eqref{eq:sing} and~\eqref{eq:comp} gives $\sum_{j\in N} x_j \le c(\{i^{\star}\}) + c(N\setminus\{i^{\star}\}) = \cstar(N) + \delta_{i^{\star}}$. But the efficiency constraint forces $\sum_j x_j = C(N)_{\text{online}} = r\cdot \cstar(N)$. Hence $r\le 1 + \delta_{i^{\star}}/\cstar(N) = r^{\ast\ast}$, contradicting the hypothesis $r>r^{\ast\ast}$.
\end{proof}

Because $\min_i \delta_i$ is typically small ($\delta_i=0$ corresponds to $\{i\},N\setminus\{i\}$ partitioning the grand tour with no overhead), $r^{\ast\ast}$ is close to $1$ and thus an operationally informative threshold. We stress that Theorem~\ref{thm:empty-core} is tight only \emph{within the single-complement mechanism}: even when no single complement $N\setminus\{i\}$ alone violates the Core, a balanced collection of complements can; this is captured by Proposition~\ref{prop:bondareva-complement} below. Intermediate-size coalitions $S$ with $|S|\in\{2,\ldots,n-2\}$ can also independently violate the Core inequality $\sum_{i\in S} x_i \le c(S)$; Section~\ref{sec:experiments} exhibits observed cases of both mechanisms (Remark~\ref{rem:near-complement} and \S\ref{ssec:cor16}).

\begin{proposition}[Balanced-complement threshold]
\label{prop:bondareva-complement}
Suppose every complement coalition is feasible, i.e., $N\setminus\{i\}\in\F$ for all $i\in N$. Define
\[
r^{\ast\ast\ast} \;:=\; \frac{1}{n-1}\cdot\frac{\sum_{i\in N} c(N\setminus\{i\})}{\cstar(N)}.
\]
If $r = C(N)_{\text{online}}/\cstar(N) > r^{\ast\ast\ast}$, then the Temporal Core is empty.
\end{proposition}

\begin{proof}
Suppose $x\in\mathrm{Core}(N,c,\F)$. For each $i\in N$, the Core constraint on $N\setminus\{i\}$ reads $\sum_{j\ne i} x_j \le c(N\setminus\{i\})$. Summing over $i\in N$ and exchanging the order of summation gives
\[
(n-1)\sum_{j\in N} x_j \;=\; \sum_{i\in N}\sum_{j\ne i} x_j \;\le\; \sum_{i\in N} c(N\setminus\{i\}).
\]
The efficiency equality $\sum_j x_j = C(N)_{\text{online}}$ yields $(n-1)\,C(N)_{\text{online}} \le \sum_i c(N\setminus\{i\})$. Dividing by $(n-1)\,\cstar(N)$ produces $r \le r^{\ast\ast\ast}$, contradicting $r>r^{\ast\ast\ast}$.
\end{proof}

The collection $\{N\setminus\{i\}\}_{i\in N}$ is balanced with weights $\lambda_i = 1/(n-1)$ (each player appears in exactly $n-1$ complements), so Proposition~\ref{prop:bondareva-complement} is the Bondareva--Shapley criterion \citep{shapley1971core} specialised to this collection. Theorem~\ref{thm:empty-core} and Proposition~\ref{prop:bondareva-complement} are incomparable in general: Theorem~\ref{thm:empty-core} depends on the minimum $\delta_i$ and fires when a single complement is cheap, whereas Proposition~\ref{prop:bondareva-complement} averages the complement costs and fires when the balanced collection as a whole is cheap. Section~\ref{sec:experiments} exhibits instances where $r > r^{\ast\ast\ast}$ but $r\le r^{\ast\ast}$, so the Core is certified empty by Proposition~\ref{prop:bondareva-complement} alone.

\begin{corollary}[Fully sequential arrivals disarm Theorem~\ref{thm:empty-core}]
\label{cor:cor51}
If arrivals are fully sequential in the sense of Lemma~\ref{lem:lemma52}, then Theorem~\ref{thm:empty-core} is vacuous: no $N\setminus\{i\}$ is feasible, so its hypothesis cannot hold.
\end{corollary}

\begin{proof}
Direct from Lemma~\ref{lem:lemma52} and the hypothesis of Theorem~\ref{thm:empty-core}. The same argument rules out Proposition~\ref{prop:bondareva-complement}'s hypothesis, which additionally requires every $N\setminus\{i\}$ to be feasible; the balanced-complement mechanism is likewise blocked under fully sequential arrivals.
\end{proof}

\begin{remark}[Near-complement mechanism (empirically observed)]
\label{rem:near-complement}
Theorem~\ref{thm:empty-core} and Proposition~\ref{prop:bondareva-complement} together cover the single-complement and balanced-complement mechanisms, i.e.\ empty Cores certifiable through coalitions of size $n-1$ alone. Beyond these, we observe main-grid instances (10 of 175 under the NN policy) where the Temporal Core is empty, both thresholds are vacuous ($r\le r^{\ast\ast}$ and $r\le r^{\ast\ast\ast}$), and the binding constraints in the Core LP mix size-$(n-1)$ and size-$(n-2)$ coalitions simultaneously---neither family alone certifies emptiness, but their joint action does. Quantitative evidence appears in Section~\ref{ssec:cor16}; Appendix~\ref{app:restricted-lp} gives the restricted-LP methodology for certifying Core emptiness at scales where $2^n$ enumeration is infeasible. A tight sufficient condition for this near-complement mechanism remains open; a natural candidate is a Bondareva--Shapley-style balanced collection mixing size-$(n-1)$ and size-$(n-2)$ coalitions (Section~\ref{sec:conclusion}).
\end{remark}

\begin{observation}[Intermediate-coalition mechanism]
\label{obs:intermediate}
In 9 of the 175 main-grid instances under the NN policy, the Temporal Core is empty despite (i) the hypothesis of Theorem~\ref{thm:empty-core} being vacuous ($r\le r^{\ast\ast}$, or no complement feasible) and (ii) the hypothesis of Proposition~\ref{prop:bondareva-complement} being vacuous. All 9 satisfy $k<n-1$ and arise in patterns $\text{B}_{\text{medium}}$ (7 cases) and E (2 cases) with realized competitive ratio $r\in[1.33,1.61]$. The binding constraints in the Core LP concentrate on intermediate-size coalitions ($2\le |S|\le n-3$) rather than on complements or near-complements. This mechanism is distinct from Theorem~\ref{thm:empty-core} and Proposition~\ref{prop:bondareva-complement} and shows that the structural safeguard of Corollary~\ref{cor:cor52}, while analytically sharp in blocking all complement-based mechanisms, does not rule out Core emptiness via intermediate coalitions when the online competitive ratio is large. An analytic sufficient condition is open (Section~\ref{sec:conclusion}).
\end{observation}

\begin{theorem}[Asymptotic tightening of $r^{\ast\ast}$]
\label{thm:asymptotic}
Assume customer coordinates $x_1,\ldots,x_n$ are drawn i.i.d.\ from the uniform distribution on $[0,L]^2$, with the depot placed at the origin $(0,0)$. Suppose the arrival pattern is such that $N\setminus\{i\}\in\F$ for at least one index $i$. Then
\[
r^{\ast\ast} \;=\; 1 + \frac{\min_{\,i:\,N\setminus\{i\}\in\F} \delta_i}{\cstar(N)}\;\xrightarrow{\text{a.s.}}\; 1 \quad \text{as } n\to\infty,
\]
with rate $r^{\ast\ast}-1 = O(n^{-1/2})$ almost surely, uniformly over the arrival patterns satisfying the hypothesis.
\end{theorem}

\begin{proof}
Fix any index $i^{0}$ with $N\setminus\{i^{0}\}\in\F$. We bound $\delta_{i^{0}}$ first, then scale by $\cstar(N)$.

\emph{Step 1: insertion and monotonicity bounds.}
Inserting the detour $\text{depot}\to x_{i^{0}}\to\text{depot}$ into any tour of $N\setminus\{i^{0}\}$ shows
\begin{equation}\label{eq:insert}
\cstar(N) \;\le\; c(N\setminus\{i^{0}\}) + c(\{i^{0}\}),
\end{equation}
whence $\cstar(N) - c(N\setminus\{i^{0}\}) \le c(\{i^{0}\}) = 2\,d(0,x_{i^{0}})$.
Conversely, removing $x_{i^{0}}$ from the optimal tour on $N$ and short-cutting its two tour neighbors yields a valid Hamiltonian cycle on $N\setminus\{i^{0}\}\cup\{0\}$ of no greater length, by the triangle inequality~\citep[\S 2.1]{steele1997probability}; hence $c(N\setminus\{i^{0}\}) \le \cstar(N)$ and $\cstar(N) - c(N\setminus\{i^{0}\}) \ge 0$. Combining:
\[
0 \;\le\; \delta_{i^{0}} \;=\; c(\{i^{0}\}) - \bigl[\,\cstar(N) - c(N\setminus\{i^{0}\})\bigr] \;\le\; c(\{i^{0}\}) \;=\; 2\,d(0,x_{i^{0}}).
\]
Since $x_{i^{0}}\in[0,L]^2$, $d(0,x_{i^{0}})\le L\sqrt{2}$ and thus $\delta_{i^{0}} \le 2L\sqrt{2} = O(1)$, independent of the arrival pattern.

\emph{Step 2: BHH lower bound on $\cstar(N)$.} By the Beardwood--Halton--Hammersley theorem \citep{beardwood1959shortest,steele1981subadditive}, the length of the optimal Hamiltonian cycle through $n$ i.i.d.\ uniform points in $[0,L]^2$ satisfies
\[
\frac{\cstar(N)}{\sqrt{n}} \;\xrightarrow{\text{a.s.}}\; \beta_{\mathrm{BHH}}\,L \quad \text{as } n\to\infty,
\]
for a universal constant $\beta_{\mathrm{BHH}}>0$. Adding a fixed depot shifts the length by $O(1)$, which is absorbed into the $o(\sqrt{n})$ term; hence $\cstar(N) = \Theta(\sqrt{n})$ a.s.

\emph{Step 3: ratio.} Combining Steps~1--2,
\[
r^{\ast\ast}-1 \;=\; \frac{\min_{\,i:\,N\setminus\{i\}\in\F}\delta_i}{\cstar(N)} \;\le\; \frac{\delta_{i^{0}}}{\cstar(N)} \;\le\; \frac{2L\sqrt{2}}{\beta_{\mathrm{BHH}}\,L\,\sqrt{n}\,(1+o(1))} \;=\; O(n^{-1/2}) \quad\text{a.s.}
\]
The numerator bound from Step~1 uses only the fixed coordinates $x_{i^{0}}\in[0,L]^2$ and not the arrival pattern, so the rate holds uniformly across patterns meeting the hypothesis. Non-negativity $r^{\ast\ast}-1\ge 0$ from Step~1 completes the claim.

\emph{Relation to the N2 convention.} The proof is stated for an arbitrary fixed coordinate square $[0,L]^2$, yielding $\delta_{i^{0}}=O(L)$ (Step~1) and $\cstar(N)=\Theta(L\sqrt{n})$ (Step~2). The parameter $L$ therefore cancels exactly in the Step~3 ratio, so the $O(n^{-1/2})$ rate is independent of the choice of $L$ as a function of $n$. Under the N2 convention $L=\sqrt{n}$ adopted in Section~\ref{ssec:design}, the numerator becomes $O(\sqrt{n})$ and the denominator $\Theta(n)$, whose ratio is the same $O(n^{-1/2})$. The theorem applies verbatim to the experimental setting.
\end{proof}

\begin{remark}[Faster decay when a depot-close complement is feasible]
\label{rem:asymptotic-tighter}
Let $i^{\star}\in\arg\min_i d(0,x_i)$ denote the customer closest to the depot. Under i.i.d.\ uniform sampling on $[0,L]^2$, the minimum of $n$ points' distances to the origin satisfies $\min_i d(0,x_i) = \Theta(n^{-1/2})$ almost surely (by standard extreme-value arguments; the radial density of uniform points near the depot satisfies $p(r)\propto r$ as $r\to 0$, whether the depot sits inside the square or at its corner, yielding the same $n^{-1/2}$ rate up to a constant). If, in addition, $N\setminus\{i^{\star}\}\in\F$---a condition automatically met when all customers share a waiting window, and more broadly whenever the arrival pattern admits the depot-closest customer's complement---Step~1 of the proof applied to $i^{\star}$ yields $\delta_{i^{\star}} \le 2\,d(0,x_{i^{\star}}) = \Theta(n^{-1/2})$ a.s., hence
\[
r^{\ast\ast}-1 \;=\; \frac{\min_i \delta_i}{\cstar(N)} \;\le\; \frac{\delta_{i^{\star}}}{\cstar(N)} \;=\; O(n^{-1}) \quad\text{a.s.,}
\]
a faster rate than the generic $O(n^{-1/2})$ of Theorem~\ref{thm:asymptotic}.

Operationally, the faster $O(n^{-1})$ rate applies primarily to arrival patterns close to the simultaneous regime, where Theorem~\ref{thm:empty-core}'s hypothesis is easily met. In sparser regimes (sequential patterns, large queue gaps) $r^{\ast\ast}$ is not defined in the first place---no $N\setminus\{i\}$ need be feasible---and the Remark is vacuous; Corollary~\ref{cor:cor52}'s structural obstruction governs those regimes.
\end{remark}

\subsection{Conditions under which the Core survives}
\label{ssec:survive}

\begin{theorem}[Inheritance of Static Core stability]
\label{thm:core-inheritance}
If $C(N)_{\text{online}} = \cstar(N)$, i.e.\ the online tour is optimal, then every static-TSG Core allocation remains in the Temporal Core; in particular $\mathrm{Core}_{\text{static}}(N,c)\subseteq\mathrm{Core}(N,c,\F)$.
\end{theorem}

\begin{proof}
Let $x\in\mathrm{Core}_{\text{static}}$. Efficiency is inherited since $\sum_i x_i = c(N) = \cstar(N) = C(N)_{\text{online}}$. For any $S\in\F\subseteq 2^N\setminus\{\emptyset\}$, the static constraint $\sum_{i\in S} x_i\le c(S)$ is already imposed on $x$, so it also holds in the temporal setting.
\end{proof}

\begin{corollary}[Structural obstruction of the complement mechanism]
\label{cor:cor52}
If the peak waiting queue satisfies $k=\max_t |U_t| < n-1$, then no coalition of the form $N\setminus\{i\}$ is feasible and Theorem~\ref{thm:empty-core} cannot fire. The only remaining \emph{a priori} mechanism for emptying the Temporal Core under $k<n-1$ is Remark~\ref{rem:remark1}, which is dominated by $r^{\ast}$ and rarely binds in practice. Intermediate-coalition violations remain logically possible; their empirical frequency is characterized in Section~\ref{sec:experiments}.
\end{corollary}

\begin{proof}
If $N\setminus\{i\}\in\F$ for some $i$, then at some instant $t$ the queue $U_t$ contains all $j\neq i$, so $|U_t|\ge n-1$ and $k\ge n-1$, contradicting $k<n-1$. The second claim follows from Theorem~\ref{thm:empty-core} and Remarks~\ref{rem:remark1}--\ref{rem:rstar}.
\end{proof}

\subsection{Discussion}
\label{ssec:discussion}

Parameters interact as follows.
\begin{itemize}[leftmargin=1.5em]
\item The \emph{tight} threshold $r^{\ast\ast}$ is sensitive only to the cheapest complement coalition: a single pair $(i,N\setminus\{i\})\in\F$ with small $\delta_i$ is enough to expose the Core.
\item The structural parameter $k$ bounds queue depth. Since $N\setminus\{i\}$-feasibility requires $|U_t|\ge n-1$, the inequality $k<n-1$ is exactly the combinatorial negation of Theorem~\ref{thm:empty-core}'s hypothesis.
\item Arrival patterns interpolate between two extremes: simultaneous ($k=n$, $\F=2^N\setminus\{\emptyset\}$, classical TSG) and fully sequential ($k\le 1$ in the steady state, Theorem~\ref{thm:empty-core} vacuous).
\end{itemize}
In the language of queueing theory, if customer inter-arrival rate $\lambda$ exceeds service throughput by enough that $k\ge n-1$ with positive probability, the complement-coalition empty-Core mechanism becomes active; otherwise it is blocked by Corollary~\ref{cor:cor52}. This provides a useful design lever: dispatchers that enforce short promise times (thereby serving orders before too many accumulate) structurally rule out the dominant empty-Core mechanism, though intermediate-coalition and near-complement violations (Remark~\ref{rem:near-complement}) remain logically possible.

\section{Computational Experiments}
\label{sec:experiments}

We report a computational study whose purpose is to (i) verify that the theoretical conditions fire exactly as predicted, (ii) quantify how much sharper $r^{\ast\ast}$ is than $r^{\ast}$, and (iii) characterize the empirical frequency of Core existence across realistic arrival regimes.

\subsection{Design}
\label{ssec:design}

\paragraph{Dimensionless framework.} Instances follow the standard dimensionless normalization of \citet{steele1997probability}: customers are sampled uniformly in $[0,L]^2$ with $L=\sqrt{n}$, fixing the customer density at one per unit area. Vehicle speed is normalized to one (unit distance per unit time), and the depot is at the origin. The system is governed by a single dimensionless load parameter
\[
  \rho \;=\; \frac{L/v}{\tau} \;=\; \frac{\text{characteristic travel time}}{\text{mean inter-arrival interval}},
\]
where $\tau$ is the mean inter-arrival interval of a given pattern. All theoretical bounds and empirical statistics depend on $(n,\rho)$ alone; the joint rescaling $(L,v,\tau)\to(\alpha L,\alpha v,\alpha\tau)$ for any $\alpha>0$ leaves all observables invariant. We verify this invariance empirically to machine precision in Appendix~\ref{app:scale-invariance}.

\paragraph{Arrival patterns.} Seven arrival patterns span the $\rho$-spectrum:
\begin{itemize}
  \item \emph{A} ($\rho=\infty$): all customers arrive simultaneously at $t=0$.
  \item \emph{$\text{B}_{\text{heavy}}$} ($\rho=4$), \emph{$\text{B}_{\text{medium}}$} ($\rho=2$), \emph{$\text{B}_{\text{light}}$} ($\rho=0.5$): sequential arrivals at uniform interval $\tau=\beta_{\mathrm{BHH}}/\rho$, where $\beta_{\mathrm{BHH}}\approx 0.7124$ is the Beardwood--Halton--Hammersley constant.
  \item \emph{C} ($\rho=2$): clustered arrivals in two angular cohorts, interleaved in space.
  \item \emph{D} ($\rho=1$): reverse-angle zig-zag order.
  \item \emph{E} ($\rho=1$): uniformly random arrivals in $[0,n\tau]$.
\end{itemize}
Patterns $\text{B}_{\text{medium}}$ and $\text{C}$ share $\rho=2$ but differ in arrival structure; Section~\ref{ssec:pattern} shows that this structural difference alone suffices to place them in qualitatively distinct regimes of Core stability, illustrating that the load parameter $\rho$ is necessary but not sufficient to characterize dynamics.

\paragraph{Computational setup.} The online tour is computed by the nearest-neighbor heuristic; Section~\ref{ssec:policy} compares this baseline with cheapest-insertion and batch re-optimization policies. The static TSP $\cstar(S)$ for each $S$ is computed exactly by dynamic programming (Held--Karp) for $|S|\le 15$ and by the LKH heuristic \citep{helsgaun2000effective} otherwise. The Temporal Nucleolus is computed by the sequential LP~\eqref{eq:lp-q} implemented with PuLP and the CBC solver.

Table~\ref{tab:design} summarizes the design: 7 patterns $\times$ 5 sizes $\times$ 5 seeds $=$ 175 instances per dispatch policy, and 3 policies yield 525 policy-instance pairs in the main study, plus an additional scale-up grid at $n\in\{20,30,50\}$ reported in Section~\ref{ssec:scaleup}.

\begin{table}[t]
\centering
\caption{Experimental design: 7 patterns $\times$ 5 sizes $\times$ 5 seeds $=$ 175 instances per dispatch policy; 3 policies yield 525 policy-instance pairs.}
\label{tab:design}
\small
\begin{tabular}{@{}ll@{}}
\toprule
Factor & Levels \\
\midrule
$n$ (coalition size) & $5,\,7,\,10,\,12,\,15$ \\
Arrival pattern & A ($\rho=\infty$), $\text{B}_{\text{heavy}}$ ($\rho=4$), $\text{B}_{\text{medium}}$ ($\rho=2$), $\text{B}_{\text{light}}$ ($\rho=0.5$), \\
                & C ($\rho=2$, clustered), D ($\rho=1$, reverse), E ($\rho=1$, random) \\
Seed & $\{7,\,42,\,99,\,123,\,256\}$ \\
Dispatch policy & Nearest-neighbor (NN), cheapest-insertion (CI), batch re-optimization (BR) \\
\bottomrule
\end{tabular}
\end{table}

\subsection{Verification of Theorem~\ref{thm:static-equiv}}

For the 25 pattern-A (simultaneous) instances we compared the Temporal Nucleolus with the static Nucleolus computed on the unrestricted coalition family $2^N\setminus\{\emptyset\}$. The two coincide coordinate-wise in every instance (maximum absolute deviation $<10^{-9}$), confirming Theorem~\ref{thm:static-equiv}.

\subsection{Verification of Theorem~\ref{thm:empty-core}}

Table~\ref{tab:thm5} is the four-mechanism contingency of predicted versus observed Core status under NN dispatch on the full 175-instance main grid.

\begin{table}[t]
\centering
\caption{Four-mechanism contingency: each of the 175 NN-policy instances classified by its empty-Core mechanism, against the observed Core status. Mechanism labels correspond to Theorem~\ref{thm:empty-core} (single-complement), Proposition~\ref{prop:bondareva-complement} (balanced-complement), Remark~\ref{rem:near-complement} (near-complement), and Observation~\ref{obs:intermediate} (intermediate). The four mechanisms together cover all 67 observed empty Cores; no false positives arise.}
\label{tab:thm5}
\small
\begin{tabular}{@{}lcc|c@{}}
\toprule
Classification & Core empty (obs.) & Core nonempty (obs.) & Total \\
\midrule
Single-complement (Theorem~\ref{thm:empty-core} fires)          & 37 &   0 &  37 \\
Balanced-complement (Proposition~\ref{prop:bondareva-complement} only) & 11 &   0 &  11 \\
Near-complement (Remark~\ref{rem:near-complement})              & 10 &   0 &  10 \\
Intermediate (Observation~\ref{obs:intermediate})               &  9 &   0 &   9 \\
No mechanism predicts emptiness                                  &  0 & 108 & 108 \\
\midrule
Total & 67 & 108 & 175 \\
\bottomrule
\end{tabular}
\end{table}

Every instance classified by any of the four mechanisms as ``empty'' is indeed empty (no false positives across 525 policy-instance pairs). Theorem~\ref{thm:empty-core} and Proposition~\ref{prop:bondareva-complement} are \emph{sufficient} conditions covering 48 of 67 NN empties; the remaining 19 split into 10 near-complement and 9 intermediate cases that are certified only by the Core LP. Figure~\ref{fig:r-rss} plots $r$ versus $r^{\ast\ast}$: instances strictly above the diagonal ($r>r^{\ast\ast}$) are the single-complement mechanism, while balanced-complement, near-complement, and intermediate cases lie at or below the diagonal and are covered by their respective certificates. The sharpness of $r^{\ast\ast}$ and $r^{\ast\ast\ast}$ relative to $r^{\ast}$ is illustrated in Figure~\ref{fig:rhist}.

\begin{figure}[t]
\centering
\includegraphics[width=0.65\textwidth]{fig2_r_vs_rstar.pdf}
\caption{Empirical test of Theorem~\ref{thm:empty-core}. The dashed line is $r=r^{\ast\ast}$; every instance strictly above the diagonal has an empty Core (red $\times$). Seven additional empty-Core instances lie at or below the diagonal, corresponding to the near-complement cases of Section~\ref{ssec:cor16} and Remark~\ref{rem:near-complement}.}
\label{fig:r-rss}
\end{figure}

\subsection{Verification of Corollary~\ref{cor:cor52} and observed intermediate-coalition mechanism}
\label{ssec:cor16}

\paragraph{Analytic sharpness of Corollary~\ref{cor:cor52}.} The analytic content of Corollary~\ref{cor:cor52} is that both complement-based mechanisms---single (Theorem~\ref{thm:empty-core}) and balanced (Proposition~\ref{prop:bondareva-complement})---are blocked whenever the peak queue satisfies $k<n-1$, because both hypotheses require at least one complement $N\setminus\{i\}$ to be feasible, and feasibility forces $|U_t|=n-1$ at some $t$. Empirically this holds without exception in the 175-instance grid: no NN (or CI, BR) instance with $k<n-1$ satisfies $r>r^{\ast\ast}$ or $r>r^{\ast\ast\ast}$ (false-positive count is zero for both thresholds across all 525 policy-instance pairs).

\paragraph{Peak-queue strata.} Partitioning the 175 main-grid instances (NN) by peak-queue regime yields three strata (Figure~\ref{fig:core-k}):
\begin{itemize}
  \item $k<n-1$: 79 instances, $70$ Core nonempty ($88.6\%$).
  \item $k=n-1$: 16 instances, $16$ Core nonempty ($100\%$).
  \item $k=n$:   80 instances, $22$ Core nonempty ($27.5\%$).
\end{itemize}
A per-pattern breakdown across $n\in\{5,7,10,12,15\}$ appears in Appendix~\ref{app:aux} (Figure~\ref{fig:core-n}).

\paragraph{Nine empty Cores at $k<n-1$: intermediate-coalition mechanism.} The 9 empty-Core instances in the $k<n-1$ regime are all certified by the intermediate-coalition mechanism of Observation~\ref{obs:intermediate}. All 9 lie in $\text{B}_{\text{medium}}$ (7 cases) or E (2 cases), have $k$ between $n-5$ and $n-2$, and realize $r\in[1.33,1.61]$. Since both Theorem~\ref{thm:empty-core} and Proposition~\ref{prop:bondareva-complement} are vacuous in the $k<n-1$ regime, the 9 cases constitute direct evidence of a third mechanism operating outside the complement-coalition framework. Table~\ref{tab:intermediate} summarises these instances; full binding-size distributions are deferred to Appendix~\ref{app:restricted-lp}.

\begin{table}[h]
\centering
\small
\caption{Nine main-grid instances with $k<n-1$ and empty Core (intermediate-coalition mechanism, Observation~\ref{obs:intermediate}). All arise under NN dispatch in patterns $\text{B}_{\text{medium}}$ or E. The ``dominant size'' column reports the coalition size at which the largest number of constraints are tight at the first-stage LP optimum.}
\label{tab:intermediate}
\begin{tabular}{@{}rllrrrrr@{}}
\toprule
$n$ & pattern & seed & $k$ & $n-1$ & $r$ & $\varepsilon^{\ast}$ & dominant size \\
\midrule
 7 & $\text{B}_{\text{medium}}$ & 256 &  5 &  6 & 1.436 & 0.028 & 4 \\
10 & $\text{B}_{\text{medium}}$ & 123 &  8 &  9 & 1.423 & 0.192 & 5 (tied with 4) \\
10 & $\text{B}_{\text{medium}}$ & 256 &  8 &  9 & 1.358 & 0.217 & 7 (tied with 6) \\
10 & E                          &  42 &  8 &  9 & 1.495 & 0.062 & 3 \\
12 & $\text{B}_{\text{medium}}$ &  99 & 10 & 11 & 1.395 & 0.744 & 6 (tied with 7) \\
12 & $\text{B}_{\text{medium}}$ & 123 & 10 & 11 & 1.487 & 0.767 & 6 \\
12 & $\text{B}_{\text{medium}}$ & 256 & 10 & 11 & 1.423 & 0.941 & 7 (tied with 8) \\
12 & E                          &  42 &  9 & 11 & 1.610 & 0.289 & 4 (tied with 5, 6) \\
15 & $\text{B}_{\text{medium}}$ & 256 & 12 & 14 & 1.326 & 0.185 & 9 (tied with 10) \\
\bottomrule
\end{tabular}
\end{table}

\paragraph{Four-way decomposition of 67 empty Cores (NN).} The 67 empty Cores decompose into four mechanisms:
\begin{itemize}
  \item Single-complement (Theorem~\ref{thm:empty-core}): \textbf{37} instances, all at $k=n$.
  \item Balanced-complement (Proposition~\ref{prop:bondareva-complement}): \textbf{11} instances, all at $k=n$.
  \item Near-complement (Remark~\ref{rem:near-complement}): \textbf{10} instances, all at $k=n$.
  \item Intermediate (Observation~\ref{obs:intermediate}): \textbf{9} instances, all at $k<n-1$.
\end{itemize}
Two regimes emerge: the saturated regime $k=n$ concentrates the three complement-based mechanisms (single + balanced + near-complement: 58 of 58 empty cases at $k=n$), while the shallow-queue regime $k<n-1$ is governed exclusively by the intermediate mechanism. These regimes are empirically disjoint in our grid.

\paragraph{Practical implication.} Corollary~\ref{cor:cor52} guarantees that keeping $k<n-1$ eliminates the complement-coalition mechanisms. In this regime, however, the intermediate mechanism of Observation~\ref{obs:intermediate} remains active and can empty the Core when the realized competitive ratio is large. Crucially, all 9 intermediate cases arise under nearest-neighbor dispatch: under cheapest-insertion and batch re-optimization, no intermediate-coalition empty Cores are observed in our grid (Section~\ref{ssec:policy}). The intermediate mechanism thus appears \emph{policy-sensitive} and can be mitigated operationally, though an analytic threshold is open.

\begin{figure}[t]
\centering
\includegraphics[width=0.7\textwidth]{fig3_core_vs_k.pdf}
\caption{Temporal Core existence rate stratified by peak-queue regime $k$. The complement-coalition mechanisms (Theorem~\ref{thm:empty-core} and Proposition~\ref{prop:bondareva-complement}) require $k\ge n-1$. Empirically, 58 of the 67 empty Cores concentrate at $k=n$ via complement / balanced-complement / near-complement mechanisms; the remaining 9 occur at $k<n-1$ via the intermediate-coalition mechanism of Observation~\ref{obs:intermediate}, all in patterns $\text{B}_{\text{medium}}$ or E.}
\label{fig:core-k}
\end{figure}

\subsection{Pattern-level analysis}
\label{ssec:pattern}

Table~\ref{tab:pattern} reports the salient statistics per arrival pattern, pooled over $n$ and seeds. Figure~\ref{fig:coal-red} visualizes the relative size of $\F$.

\begin{table}[t]
\centering
\caption{Pattern-level aggregate statistics (NN policy, $5\times 5=25$ instances per pattern). $\bar r^{\ast}$, being a function of coordinates only ($r^{\ast} = \sum_i c(\{i\})/\cstar(N)$), is identical across all patterns at $4.351$ and reported as an overall value rather than per-pattern. Em-dashes in $\bar r^{\ast\ast}$ (resp.\ $\bar r^{\ast\ast\ast}$) indicate patterns where Theorem~\ref{thm:empty-core} (resp.\ Proposition~\ref{prop:bondareva-complement}) is never applicable.}
\label{tab:pattern}
\small
\begin{tabular}{@{}lcccccccc@{}}
\toprule
Pattern & $\rho$ & Core rate & $\bar r$ (sd) & $\bar r^{\ast\ast}$ (sd) & $\bar r^{\ast\ast\ast}$ & $\overline{|\F|/(2^n-1)}$ & Thm 11 app/fires & Prop 12 app/fires \\
\midrule
A                           & $\infty$ & 0.36 & 1.125 (0.102) & 1.181 (0.059) & 1.092 & 1.000 & 25 / 9  & 25 / 15 \\
$\text{B}_{\text{heavy}}$   & 4.0      & 0.20 & 1.233 (0.127) & 1.187 (0.061) & 1.094 & 0.985 & 25 / 15 & 24 / 20 \\
$\text{B}_{\text{medium}}$  & 2.0      & 0.56 & 1.273 (0.135) & 1.305 (0.148) & 1.150 & 0.590 & 13 / 3  &  7 /  4 \\
$\text{B}_{\text{light}}$   & 0.5      & 1.00 & 1.690 (0.210) &      ---      &  ---  & 0.077 &  0 / 0  &  0 /  0 \\
C                           & 2.0      & 0.28 & 1.197 (0.149) & 1.181 (0.059) & 1.087 & 0.985 & 25 / 10 & 24 / 18 \\
D                           & 1.0      & 1.00 & 1.332 (0.121) & 1.398 (0.217) &  ---  & 0.196 &  3 / 0  &  0 /  0 \\
E                           & 1.0      & 0.92 & 1.337 (0.169) & 1.505 (0.139) &  ---  & 0.273 &  5 / 0  &  0 /  0 \\
\midrule
Overall                     & ---      & 0.62 & 1.312 (0.224) & 1.223 (0.121) & 1.096 & 0.587 & 96 / 37 & 80 / 57 \\
\bottomrule
\end{tabular}
\end{table}

Three patterns stand out. Simultaneous arrivals (A, $\rho=\infty$) yield the lowest Core-existence rate ($36\%$), consistent with $\F=2^N\setminus\{\emptyset\}$ maximizing the number of Core constraints. Light-load and fully-dispersed sequential patterns ($\text{B}_{\text{light}}$, D, E) leave $|\F|$ at only $6$--$15\%$ of $2^n-1$ and render Theorem~\ref{thm:empty-core} inapplicable ($0/25$ in each). The clustered pattern (C) and the heavy-load $\text{B}_{\text{heavy}}$ both have $\rho\ge 2$ and yield $\overline{|\F|/(2^n-1)}\approx 0.5$, placing them in the saturated regime where Theorem~\ref{thm:empty-core} applies to nearly every instance but is binding in only $12\%$ (C) and $8\%$ ($\text{B}_{\text{heavy}}$) of them.

A secondary observation: the mean competitive ratio $\bar r$ rises sharply in the heavily sequential regimes ($\text{B}_{\text{light}}$: $1.69$, E: $1.34$, D: $1.33$) because the online tour takes detours to reach customers as they arrive while $\cstar(N)$ reflects the offline-optimal route. In patterns $\text{B}_{\text{light}}$ and D the queue stays shallow throughout, and Corollary~\ref{cor:cor52} blocks every complement-based mechanism; the Core remains nonempty in all 25 instances for each pattern. In pattern E the same obstruction blocks complement mechanisms, yet 2 of the 25 instances still exhibit empty Cores via the intermediate-coalition mechanism of Observation~\ref{obs:intermediate}---see \S\ref{ssec:cor16}. The operational content of Corollary~\ref{cor:cor52} is therefore precise: a shallow queue rules out complement and balanced-complement mechanisms, but does not by itself prevent intermediate-coalition violations at high $\bar r$.

\paragraph{Load parameter $\rho$ alone does not determine Core stability.} Patterns $\text{B}_{\text{medium}}$ and C share the same load parameter $\rho=2$ yet exhibit qualitatively different behavior. $\text{B}_{\text{medium}}$ (uniform sequential arrivals) is complement-mechanism-free in most instances (Theorem~\ref{thm:empty-core} is applicable in only $13/25$ and fires in $3/25$), yet 7 further instances are empty through the intermediate mechanism, giving an overall Core rate of $0.56$. C (two angularly-interleaved batches) places almost every instance ($24/25$) into the Theorem~\ref{thm:empty-core}-applicable regime, where single-complement violations fire in $10/25$, producing a Core rate of $0.28$. The structural difference in the arrival geometry alone accounts for this gap. This confirms that $\rho$, while necessary for placing a pattern on the saturation spectrum, is not sufficient to characterize Core stability: the spatial-temporal structure of the arrival process and the dispatch-policy dynamics both matter independently.


\begin{figure}[t]
\centering
\includegraphics[width=0.85\textwidth]{fig4_coalition_reduction.pdf}
\caption{Relative size of the feasible coalition family, $|\F|/(2^n-1)$.}
\label{fig:coal-red}
\end{figure}

\subsection{\texorpdfstring{Sharpness of $r^{\ast\ast}$ versus $r^{\ast}$}{Sharpness of r** versus r*}}

Figure~\ref{fig:rhist} contrasts the distributions of the single-complement threshold $r^{\ast\ast}$ (Theorem~\ref{thm:empty-core}), the balanced-complement threshold $r^{\ast\ast\ast}$ (Proposition~\ref{prop:bondareva-complement}), and the classical bound $r^{\ast}$ (Remark~\ref{rem:rstar}). Empirically, on the NN grid $\bar r^{\ast\ast}=1.223$ and $\bar r^{\ast\ast\ast}=1.096$ versus $\bar r^{\ast}=4.351$; the sharpness factors are $r^{\ast}/r^{\ast\ast} \approx 3.56$ and $r^{\ast}/r^{\ast\ast\ast} \approx 3.97$. Because the realized ratio $\bar r = 1.312$ is close to $r^{\ast\ast}$ and $r^{\ast\ast\ast}$ but far below $r^{\ast}$, Theorem~\ref{thm:empty-core} and Proposition~\ref{prop:bondareva-complement} are operationally informative whereas Remark~\ref{rem:rstar} almost never fires. The classical bound grows with $n$ because $r^{\ast} = \sum_i c(\{i\})/\cstar(N)$ has a numerator scaling linearly in $n$ while $\cstar(N)$ scales like $\sqrt{n}$ on uniform planar instances; $r^{\ast\ast}$ and $r^{\ast\ast\ast}$, in contrast, are pinned near $1$ because complement coalitions $N\setminus\{i\}$ nearly partition the grand tour. This explains why the gap between the two widens with $n$ and quantifies the practical advance of the temporal analysis. Note that $r^{\ast\ast\ast}$ averages across complements while $r^{\ast\ast}$ takes the minimum of $\delta_i$; the two thresholds catch different instances---of the 57 Proposition~\ref{prop:bondareva-complement} fires under NN, 20 are outside the Theorem~\ref{thm:empty-core} single-complement regime.

A secondary point worth noting: $r^{\ast}$ depends only on coordinates---it equals $\sum_i c(\{i\})/\cstar(N)$, and neither the numerator nor the denominator sees arrivals---so under our uniform-square sampling $\bar r^{\ast}$ has essentially no variance across patterns. The operational shortfall of $r^{\ast}$ is therefore structural: it is blind to the online arrival process altogether.

\begin{figure}[t]
\centering
\includegraphics[width=0.8\textwidth]{fig5_rstar_vs_rss.pdf}
\caption{Histograms of $r^{\ast\ast}$, $r^{\ast\ast\ast}$, and $r^{\ast}$ over all NN instances. $r^{\ast\ast}$ and $r^{\ast\ast\ast}$ concentrate near the empirical $r$ (making them informative Core-emptiness thresholds via Theorem~\ref{thm:empty-core} and Proposition~\ref{prop:bondareva-complement}); $r^{\ast}$ lies far to the right and rarely binds.}
\label{fig:rhist}
\end{figure}

\subsection{Robustness to dispatch policy}
\label{ssec:policy}

We rerun the 175-instance grid under three dispatch policies: nearest-neighbor (NN, the baseline used in the preceding subsections), cheapest-insertion (CI), and batch re-optimization (BR, which solves the shortest Hamiltonian path from the vehicle's current position through the waiting queue $U_t$ back to the depot by Held--Karp at each dispatch event). Table~\ref{tab:policy} summarizes the aggregate outcomes on all $175\times 3 = 525$ policy-instance pairs.

\begin{table}[t]
\centering
\caption{Dispatch-policy robustness over the 175-instance grid. $\bar r$ is the mean competitive ratio; the last columns report how often each analytic mechanism fires and how many intermediate-coalition empty-Core cases (Observation~\ref{obs:intermediate}) arise.}
\label{tab:policy}
\small
\begin{tabular}{@{}lccccccc@{}}
\toprule
Policy & $\bar r$ & Core rate & Empty & Thm~\ref{thm:empty-core} app/fires & Prop~\ref{prop:bondareva-complement} app/fires & Obs~\ref{obs:intermediate} \\
\midrule
Nearest-neighbor      & 1.312 & 0.617 & 67 & 96/37 & 80/57 & 9 \\
Cheapest-insertion    & 1.223 & 0.823 & 31 & 97/17 & 80/30 & 0 \\
Batch re-optimization & 1.201 & 0.857 & 25 & 97/ 7 & 80/25 & 0 \\
\bottomrule
\end{tabular}
\end{table}

Three observations stand out. First, the mean competitive ratio decreases monotonically from $1.312$ (NN) to $1.223$ (CI) to $1.201$ (BR), reflecting the increasing quality of the online plan. Second, the empirical Core-existence rate increases correspondingly: NN yields 67 empty Cores ($38\%$), CI 31 ($18\%$), and BR 25 ($14\%$). Third, and central to our theoretical claims, \emph{every} instance in which the Theorem~\ref{thm:empty-core} hypothesis $r>r^{\ast\ast}$ holds is observed empty (37/37 for NN, 17/17 for CI, 7/7 for BR), and \emph{every} instance in which Proposition~\ref{prop:bondareva-complement}'s hypothesis $r>r^{\ast\ast\ast}$ holds is likewise observed empty (57/57, 30/30, 25/25), with no false positives under any policy. The Corollary~\ref{cor:cor52} safeguard holds analytically without exception across all three policies: no instance with $k<n-1$ has $r>r^{\ast\ast}$ or $r>r^{\ast\ast\ast}$. Intermediate-coalition empty Cores (Observation~\ref{obs:intermediate}) appear exclusively under NN; CI and BR produce \emph{zero} intermediate cases in the 175-instance grid, corroborating the policy-sensitivity of that mechanism.

Applicability counts are nearly identical across policies: Theorem~\ref{thm:empty-core} is applicable in $96$--$97$ of the $175$ instances, and Proposition~\ref{prop:bondareva-complement} in $80$ (the applicability gap is because Proposition~\ref{prop:bondareva-complement} requires \emph{every} complement $N\setminus\{i\}$ to be feasible whereas Theorem~\ref{thm:empty-core} requires only one). Applicability is mostly insensitive to policy because the feasibility of a specific $N\setminus\{i\}$ depends on the overlap of service-time windows, which is only weakly affected by dispatch order for our arrival patterns. The thresholds $r^{\ast\ast}$ and $r^{\ast\ast\ast}$ themselves are policy-independent as functions of TSP costs and arrivals (not service times), so only the realized $r$ varies with policy; $r>r^{\ast\ast}$ and $r>r^{\ast\ast\ast}$ both become rarer as the policy improves, which is exactly what Table~\ref{tab:policy} shows.

\subsection{Scale-up to larger \texorpdfstring{$n$}{n}}
\label{ssec:scaleup}

Our main study at $n\le 15$ is capped by the $O(n^2 2^n)$ cost of Held--Karp. To test whether the theoretical and empirical findings of the preceding subsections persist at larger scales, we extend the study to $n\in\{20,30,50\}$ for three representative patterns---A ($\rho=\infty$, simultaneous), $\text{B}_{\text{medium}}$ ($\rho=2$, sequential), and C ($\rho=2$, clustered)---spanning the Core-fragile, Core-safe, and intermediate regimes. Each cell uses 5 seeds, giving 45 scale-up instances with NN dispatch.

\paragraph{Solver and calibration.} For $n\le 20$ we use exact Held--Karp (pure distance); for $n>20$ we use the LKH heuristic of \citet{helsgaun2000effective} via the \texttt{elkai} Python wrapper. Calibration against Held--Karp on 15 instances at $n=15$ yielded a relative error of $0.000\%$ in every case, so at this scale LKH is effectively exact for our pure-distance $c^{\ast}(S)$.

\begin{table}[t]
\centering
\caption{Scale-up aggregate statistics (5 seeds each). Em-dashes in $\bar r^{\ast\ast}$ indicate no feasible complement coalition (Theorem~\ref{thm:empty-core} vacuous). For pattern A with $n\ge 30$ the Core LP is intractable to enumerate, since $|\F|=2^n-1$; emptiness there is certified a priori by Theorem~\ref{thm:empty-core} for the 4 of 5 seeds at $n=30$ and 5 of 5 at $n=50$ where $r>r^{\ast\ast}$, and by a restricted Core LP (Appendix~\ref{app:restricted-lp}) for the one seed at $n=30$ where $r\le r^{\ast\ast}$.}
\label{tab:scaleup}
\small
\begin{tabular}{@{}lccccc@{}}
\toprule
Pattern & $n$ & Core rate & $\bar r$ & $\bar r^{\ast\ast}$ & $\bar k$ \\
\midrule
A                          & 20 & 0 / 5$^{\ddagger}$       & 1.167 & 1.119 & 20.0 \\
A                          & 30 & 0 / 5$^{\dagger\ddagger}$ & 1.187 & 1.096 & 30.0 \\
A                          & 50 & 0 / 5$^{\dagger}$         & 1.211 & 1.068 & 50.0 \\
\midrule
$\text{B}_{\text{medium}}$ & 20 & 5 / 5 & 1.330 & --- & 15.4 \\
$\text{B}_{\text{medium}}$ & 30 & 5 / 5 & 1.309 & --- & 23.2 \\
$\text{B}_{\text{medium}}$ & 50 & 5 / 5 & 1.387 & --- & 37.6 \\
\midrule
C                          & 20 & 0 / 5 & 1.335 & 1.119 & 20.0 \\
C                          & 30 & 0 / 5 & 1.282 & 1.096 & 30.0 \\
C                          & 50 & 0 / 5 & 1.282 & 1.068 & 50.0 \\
\bottomrule
\end{tabular}

{\footnotesize $^{\dagger}$ Core LP not enumerated at this scale; Theorem~\ref{thm:empty-core} fires in $4/5$ seeds at $n=30$ and $5/5$ at $n=50$; the one non-firing seed at $n=30$ (seed 123) is certified by a restricted Core LP (Appendix~\ref{app:restricted-lp}). $^{\ddagger}$ At $n=20$, Theorem~\ref{thm:empty-core} is vacuous ($r\le r^{\ast\ast}$) at seeds 99 and 123 and fires at seeds 7, 42, 256; at $n=30$, it is vacuous at seed 123 only. Core emptiness for the vacuous seeds is certified by the near-complement mechanism (Remark~\ref{rem:near-complement}, Section~\ref{ssec:cor16}).}
\end{table}

Four findings emerge.

\paragraph{Finding 1: $\bar r^{\ast\ast}$ decays with $n$, confirming Theorem~\ref{thm:asymptotic}.} For pattern A, $\bar r^{\ast\ast}$ drops from $1.181$ at $n=15$ (Table~\ref{tab:pattern}) to $1.119$, $1.096$, $1.068$ at $n=20, 30, 50$, respectively. Pattern C exhibits the same decay ($\text{B}_{\text{medium}}$ has no feasible complement at any scale, so $\bar r^{\ast\ast}$ is undefined). The observed rate of decrease is consistent with the $O(n^{-1/2})$ bound of Theorem~\ref{thm:asymptotic} (and within the $O(n^{-1})$ improvement of Remark~\ref{rem:asymptotic-tighter} for the simultaneous-arrival case A).

\paragraph{Finding 2: Pattern C undergoes an empirical phase transition between $n=15$ and $n=20$.} At $n=15$ the clustered pattern had an $88\%$ Core-existence rate with $\bar r = 1.199 < \bar r^{\ast\ast} = 1.363$ (Table~\ref{tab:pattern}). By $n=20$ the two quantities have crossed: $\bar r = 1.335 > \bar r^{\ast\ast} = 1.119$, and all 5 test instances have an empty Core; the same holds at $n=30, 50$. Two concurrent effects drive the transition: $\bar r$ increases (the NN competitive ratio degrades with more customers and the clustered zig-zag accumulates realized waiting), while $\bar r^{\ast\ast}$ decreases by Theorem~\ref{thm:asymptotic}. Clustered patterns are thus immune to the complement-coalition mechanism only at modest $n$; the mechanism activates as $n$ grows.

\paragraph{Finding 3: Pattern $\text{B}_{\text{medium}}$ avoids the complement mechanism at every scale.} Across $n\in\{20,30,50\}$, $\text{B}_{\text{medium}}$ exhibits $\bar r \in [1.31,1.39]$ yet every instance retains a nonempty Core in our experiments. Structurally, $\text{B}_{\text{medium}}$'s sparse sequential arrivals keep $\bar k$ strictly below $n-1$ at every scale ($\bar k = 15.4, 23.2, 37.6$ at $n=20, 30, 50$), so every complement coalition $N\setminus\{i\}$ is infeasible, Theorem~\ref{thm:empty-core} is vacuous, and Corollary~\ref{cor:cor52}'s structural obstruction applies. While intermediate-coalition violations remain logically possible, none occurred in our grid. This validates the ``small-$k$ blocks the complement mechanism'' principle at scale: empirical Core stability for $\text{B}_{\text{medium}}$ does not require low $r$---a shallow queue suffices to rule out the dominant empty-Core mechanism.

\paragraph{Finding 4: Near-complement mechanism at seed 123.} At seed 123 under Pattern~A, Theorem~\ref{thm:empty-core} is vacuous at $n=20$ and $n=30$ ($r=1.077<r^{\ast\ast}=1.112$ and $r=1.038<r^{\ast\ast}=1.088$ respectively) but fires at $n=50$ ($r=1.306>r^{\ast\ast}=1.081$), consistent with the $O(n^{-1/2})$ decay of $\bar r^{\ast\ast}$. Nevertheless, a restricted Core LP over $10{,}000$ sampled coalitions certifies Core emptiness at all three scales, with binding constraints spanning both complement (size $n-1$) and near-complement (size $n-2$) coalitions simultaneously in each case. By contrast, a clean Theorem-\ref{thm:empty-core}-fires instance (Pattern A, $n=30$, seed 42) binds only at size $n-1$ (Section~\ref{ssec:cor16}, Appendix~\ref{app:restricted-lp}). The persistence of size-$(n-2)$ binding across scale---even where Theorem~\ref{thm:empty-core} is applicable---suggests that the near-complement mechanism is a stable structural feature of this configuration rather than a small-$n$ artifact.

\paragraph{Intractability of pattern A at $n\ge 30$.} For pattern A, $|\F| = 2^n-1$, so the Core LP has $\sim\!10^9$ and $\sim\!10^{15}$ constraints at $n=30$ and $n=50$ respectively, well beyond enumeration. Theorem~\ref{thm:empty-core} nevertheless certifies Core emptiness \emph{a priori} in $12$ of the $15$ seed--size combinations: every complement coalition is feasible, so $r^{\ast\ast}$ is well-defined, and numerically $r>r^{\ast\ast}$ holds in $3/5$ seeds at $n=20$, $4/5$ at $n=30$, and $5/5$ at $n=50$ (Table~\ref{tab:scaleup}). The three exceptions---seeds 99 and 123 at $n=20$, and seed 123 at $n=30$---have $r\le r^{\ast\ast}$; for these the restricted Core LP of Appendix~\ref{app:restricted-lp} certifies emptiness (Finding~4 above). The intractability of pattern A therefore itself \emph{demonstrates} the $k<n-1$ versus $k=n$ divide identified in Corollary~\ref{cor:cor52}: the regime where the Core LP becomes infeasible to enumerate is precisely the regime where Theorem~\ref{thm:empty-core}---supplemented by the near-complement certificate at the few exceptional seeds---can certify emptiness without enumeration.

\paragraph{Summary.} The qualitative findings of the earlier subsections persist up to $n=50$, with one refinement: the tail of Core-empty clustered instances at $n=15$ expands to near-certain Core emptiness by $n=20$ because $r^{\ast\ast}$ tightens with $n$ (Theorem~\ref{thm:asymptotic}). Corollary~\ref{cor:cor52}'s safeguard, in contrast, only strengthens with $n$: across every scale-up cell, whenever $\bar k<n-1$ the Core is nonempty irrespective of the realized competitive ratio (observed $\bar r$ as high as $1.39$ at $n=50$ for $\text{B}_{\text{medium}}$).

\section{Conclusion}
\label{sec:conclusion}

We introduced the Online Traveling Salesman Game, in which coalitions are restricted to those realizable under the online arrival process, and defined the Temporal Nucleolus as the lexicographic excess minimizer over this restricted family. We showed that the classical (static) TSG is recovered at the boundary case of simultaneous arrivals, that the Temporal Core is empty whenever the competitive ratio exceeds a complement-mechanism threshold $r^{\ast\ast}$, and that a bounded waiting queue ($k<n-1$) structurally blocks the complement-coalition empty-Core mechanism (Corollary~\ref{cor:cor52}), with empirical Core nonemptiness in $79$ of $79$ such instances. Numerical experiments over 175 instances confirm every theoretical prediction and show that this threshold is sharper than the classical bound by a factor of $3.21$, making it the first operationally informative threshold for cost-sharing stability in online routing \emph{within the complement-coalition mechanism}; the broader landscape includes near-complement effects (Section~\ref{ssec:cor16}, Remark~\ref{rem:near-complement}, Appendix~\ref{app:restricted-lp}) that warrant further study.

\paragraph{Platform-design implication.} From a platform-design perspective, our results offer suggestive evidence that enforcing short service-time promises---equivalently, keeping the peak queue depth below $n-1$---structurally rules out the complement-coalition empty-Core mechanism. Near-complement violations remain logically possible in the saturated regime and are empirically observed in a minority of instances (Section~\ref{ssec:cor16}).

\paragraph{Limitations and future work.} Four limitations merit explicit discussion.

\emph{(i) Dispatch policy.} Theorems~\ref{thm:empty-core} and~\ref{thm:core-inheritance} and Corollary~\ref{cor:cor52} use only the realized tuple $(r, \F, k)$ and are therefore policy-independent in their statements. Section~\ref{ssec:policy} extends the empirical study to cheapest-insertion and batch re-optimization policies: the qualitative conclusions---Theorem~\ref{thm:empty-core}'s direction of implication and Corollary~\ref{cor:cor52}'s structural obstruction---are preserved, while the absolute value of $\bar r$ varies by policy as expected. Extending to learning-based or forecasting-aware dispatch policies remains open.

\emph{(ii) Instance scale.} Section~\ref{ssec:scaleup} extends the study to $n\in\{20,30,50\}$ for patterns A, $\text{B}_{\text{medium}}$, and C, using LKH \citep{helsgaun2000effective} for $c^{\ast}(S)$ at $n>20$ with calibration error $\eta=0.000\%$ on $n=15$ pure-distance benchmarks. Four findings emerge: $\bar r^{\ast\ast}$ decays as $O(n^{-1/2})$ (Theorem~\ref{thm:asymptotic}), pattern C undergoes a phase transition between $n=15$ and $n=20$ as $\bar r$ overtakes $\bar r^{\ast\ast}$, pattern $\text{B}_{\text{medium}}$ avoids the complement-coalition empty-Core mechanism at every scale as its queue stays shallow ($\bar k<n-1$ throughout), with Core nonemptiness observed in every instance---a direct validation of Corollary~\ref{cor:cor52}, and the near-complement mechanism (Remark~\ref{rem:near-complement}) is observed at seed 123 across $n\in\{20,30,50\}$ even where Theorem~\ref{thm:empty-core} fires. Pattern A is Core-LP intractable for $n\ge 30$ because $|\F| = 2^n - 1$ there; emptiness is nevertheless certified \emph{a priori} by Theorem~\ref{thm:empty-core} (12 of 15 seed--size combinations) and by a restricted Core LP (Appendix~\ref{app:restricted-lp}) for the three exceptions (seeds 99 and 123 at $n=20$, and seed 123 at $n=30$). Extending to $n\ge 100$, to non-uniform geometries, and to explicit heavy-tailed queueing models remains for future work.

\emph{(iii) Near-complement mechanism.} Remark~\ref{rem:near-complement} documents 10 main-grid NN instances where Core emptiness is certified only by the combined effect of $(n-1)$- and $(n-2)$-size binding constraints, with neither Theorem~\ref{thm:empty-core} nor Proposition~\ref{prop:bondareva-complement} firing. A tight sufficient condition is open; we conjecture it involves balanced collections (in the sense of \citealp{shapley1971core}) mixing size-$(n-1)$ and size-$(n-2)$ coalitions, extending the single- and balanced-complement arguments of Theorem~\ref{thm:empty-core} and Proposition~\ref{prop:bondareva-complement}.

\emph{(iv) Intermediate-coalition mechanism.} Observation~\ref{obs:intermediate}'s 9 empirical cases call for an analytic sufficient condition. Two avenues suggest themselves: (a) balanced collections of intermediate-size coalitions in the sense of \citet{shapley1971core} with weights adapted to the waiting-window structure of $\F$, and (b) a queueing-theoretic characterization of when sequential arrivals under high-$r$ policies accumulate enough intermediate coalitions to block the Core. Both are deferred to future work. Our empirical evidence suggests that the mechanism is policy-sensitive: observed only under nearest-neighbor dispatch in our grid, absent under cheapest-insertion and batch re-optimization, and concentrated in patterns $\text{B}_{\text{medium}}$ and E at $n\in\{7,10,12,15\}$.

\emph{Broader directions.} Beyond these limitations, natural extensions include the Online Vehicle Routing Game with multiple vehicles and capacity constraints, and an A3-style consistency axiom so that allocations computed across rolling-horizon windows agree on overlapping coalitions.

\section*{CRediT author statement}
\textbf{Seyun Jeong:} Conceptualization, Methodology, Formal analysis, Software, Writing --- original draft. \textbf{Hyunchul Tae:} Conceptualization, Methodology, Supervision, Writing --- review \& editing.

\section*{Acknowledgments}
The authors thank collaborators for helpful discussions on earlier drafts.

\section*{Declarations of interest}
The authors declare no competing financial interests or personal relationships that could have appeared to influence the work reported in this paper.

\section*{Data and code availability}
All code, random seeds, and raw CSV results used in the computational study (Sections~\ref{sec:experiments} and~\ref{ssec:scaleup}) are available at \url{https://github.com/sage-Tae/online-tsg}. The repository contains the Held--Karp and LKH implementations, the dispatch-policy simulators for nearest-neighbor, cheapest-insertion, and batch re-optimization, and the Temporal Nucleolus LP solver, along with scripts reproducing all tables and figures in this paper.

\bibliographystyle{plainnat}
\bibliography{references}

\appendix
\section{Auxiliary plots}
\label{app:aux}

Figure~\ref{fig:core-n} breaks down the aggregate Core-existence rates of Section~\ref{sec:experiments} by arrival pattern and coalition size $n$. Patterns $\text{B}_{\text{medium}}$, $\text{B}_{\text{light}}$, D, and E retain $100\%$ Core existence across all $n\in\{5,7,10,12,15\}$. Pattern $\text{B}_{\text{heavy}}$ is at $100\%$ for $n\le 10$ and drops to $80\%$ at $n\in\{12,15\}$; the clustered pattern C dips to $80\%$ at $n=10$ and $60\%$ at $n=12$ before recovering to $100\%$ at $n=15$. Simultaneous arrivals (A) collapse from $60\%$ at $n=5$ to $0\%$ at $n=15$. Figure~\ref{fig:core-k} in the main text summarizes this by peak-queue regime $k$.

\begin{figure}[h]
\centering
\includegraphics[width=0.75\textwidth]{fig3_core_vs_n.pdf}
\caption{Core-existence rate versus coalition size $n$, broken down by arrival pattern. Patterns $\text{B}_{\text{medium}}$, $\text{B}_{\text{light}}$, D, and E overlap at $1.00$ across the grid. $\text{B}_{\text{heavy}}$ remains at $1.00$ for $n\le 10$ and drops to $0.80$ at $n\in\{12,15\}$; C dips to $0.80$--$0.60$ at $n\in\{10,12\}$ before recovering; simultaneous arrivals (A) collapse to $0\%$ at $n=15$.}
\label{fig:core-n}
\end{figure}

\section{Scale invariance: empirical verification}
\label{app:scale-invariance}

This appendix verifies empirically the scale invariance stated in Section~\ref{ssec:design}: that all dimensionless observables are invariant under the joint rescaling $(L, v, \tau) \to (\alpha L, \alpha v, \alpha \tau)$ for any $\alpha > 0$. Since we fix $v = 1$ throughout, this reduces to the simultaneous scaling $(L, \tau) \to (\alpha L, \alpha \tau)$ with the ratio $L/\tau$ (hence the load parameter $\rho$) preserved.

\paragraph{Experimental setup.} Fixing seed $=42$, we generate instances at $n\in\{10,20\}$ and patterns $\{\text{B}_{\text{heavy}},\text{B}_{\text{medium}}, C\}$ under four scale factors $\alpha\in\{0.5,1.0,2.0,5.0\}$. For each configuration, coordinates are sampled from $[0,\alpha\sqrt{n}]^2$ and arrival intervals are scaled from the pattern-prescribed base $\tau_0=\beta_{\mathrm{BHH}}/\rho$ to $\alpha\tau_0$. The nearest-neighbor dispatch policy runs on all 24 instances. Statistics $r$, $r^{\ast\ast}$, and $k$ are recorded.

\paragraph{Results.} Table~\ref{tab:scale-invariance} reports the full 24 instances. For each of the six $(n,\text{pattern})$ configurations, all four $\alpha$ values yield \emph{bit-identical} statistics: the relative standard deviation of $r$ across $\alpha$ is $0$ to machine precision, and $k$ and $r^{\ast\ast}$ are unchanged across $\alpha$ in every configuration. Only the dimensional quantities $L$ and $\tau$ scale linearly with $\alpha$ as expected.

\begin{table}[htbp]
\centering
\footnotesize
\caption{Scale invariance verification (NN policy, seed $=42$). Each row corresponds to one $(n,\text{pattern},\alpha)$ configuration; 24 rows total. Within each $(n,\text{pattern})$ block, the four $\alpha$ values yield identical $r$, $r^{\ast\ast}$, and $k$, confirming joint-rescaling invariance to machine precision. Em-dashes in $r^{\ast\ast}$ indicate configurations where Theorem~\ref{thm:empty-core} is inapplicable (no complement coalition feasible).}
\label{tab:scale-invariance}
\begin{tabular}{@{}cccccccc@{}}
\toprule
$n$ & Pattern & $\alpha$ & $L$ & $\tau$ & $r$ & $r^{\ast\ast}$ & $k$ \\
\midrule
10 & $\text{B}_{\text{heavy}}$ & 0.5 & 1.5811 & 0.0891 & 1.4225 & 1.1324 & 10 \\
10 & $\text{B}_{\text{heavy}}$ & 1.0 & 3.1623 & 0.1781 & 1.4225 & 1.1324 & 10 \\
10 & $\text{B}_{\text{heavy}}$ & 2.0 & 6.3246 & 0.3562 & 1.4225 & 1.1324 & 10 \\
10 & $\text{B}_{\text{heavy}}$ & 5.0 & 15.8114 & 0.8905 & 1.4225 & 1.1324 & 10 \\
\midrule
10 & $\text{B}_{\text{medium}}$ & 0.5 & 1.5811 & 0.1781 & 1.4225 & 1.1324 & 10 \\
10 & $\text{B}_{\text{medium}}$ & 1.0 & 3.1623 & 0.3562 & 1.4225 & 1.1324 & 10 \\
10 & $\text{B}_{\text{medium}}$ & 2.0 & 6.3246 & 0.7124 & 1.4225 & 1.1324 & 10 \\
10 & $\text{B}_{\text{medium}}$ & 5.0 & 15.8114 & 1.7810 & 1.4225 & 1.1324 & 10 \\
\midrule
10 & C & 0.5 & 1.5811 & 0.1781 & 1.2531 & 1.1324 & 10 \\
10 & C & 1.0 & 3.1623 & 0.3562 & 1.2531 & 1.1324 & 10 \\
10 & C & 2.0 & 6.3246 & 0.7124 & 1.2531 & 1.1324 & 10 \\
10 & C & 5.0 & 15.8114 & 1.7810 & 1.2531 & 1.1324 & 10 \\
\midrule
20 & $\text{B}_{\text{heavy}}$ & 0.5 & 2.2361 & 0.0891 & 1.4896 & 1.1039 & 20 \\
20 & $\text{B}_{\text{heavy}}$ & 1.0 & 4.4721 & 0.1781 & 1.4896 & 1.1039 & 20 \\
20 & $\text{B}_{\text{heavy}}$ & 2.0 & 8.9443 & 0.3562 & 1.4896 & 1.1039 & 20 \\
20 & $\text{B}_{\text{heavy}}$ & 5.0 & 22.3607 & 0.8905 & 1.4896 & 1.1039 & 20 \\
\midrule
20 & $\text{B}_{\text{medium}}$ & 0.5 & 2.2361 & 0.1781 & 1.4896 & --- & 17 \\
20 & $\text{B}_{\text{medium}}$ & 1.0 & 4.4721 & 0.3562 & 1.4896 & --- & 17 \\
20 & $\text{B}_{\text{medium}}$ & 2.0 & 8.9443 & 0.7124 & 1.4896 & --- & 17 \\
20 & $\text{B}_{\text{medium}}$ & 5.0 & 22.3607 & 1.7810 & 1.4896 & --- & 17 \\
\midrule
20 & C & 0.5 & 2.2361 & 0.1781 & 1.3469 & 1.1039 & 20 \\
20 & C & 1.0 & 4.4721 & 0.3562 & 1.3469 & 1.1039 & 20 \\
20 & C & 2.0 & 8.9443 & 0.7124 & 1.3469 & 1.1039 & 20 \\
20 & C & 5.0 & 22.3607 & 1.7810 & 1.3469 & 1.1039 & 20 \\
\bottomrule
\end{tabular}
\end{table}

\paragraph{Interpretation.} The absence of variation in $r$, $r^{\ast\ast}$, and $k$ across $\alpha$ confirms that the online arrival process, the NN tour construction, and the Core analysis all depend only on the dimensionless parameters $(n,\rho)$ and are insensitive to the absolute scale of $(L,\tau)$. This justifies the specific choice $L=\sqrt{n}$ adopted in Section~\ref{ssec:design} without loss of generality.

\section{Restricted Core LP and mechanism classification}
\label{app:restricted-lp}

This appendix (i) documents the restricted Core LP methodology used for scale-up instances where $2^n-1$ enumeration is infeasible, and (ii) gives the full binding-size distributions for the 10 near-complement and 9 intermediate main-grid cases under NN.

\paragraph{Motivation (scale-up).} The full Core LP requires one inequality per feasible coalition $S\in\mathcal{F}$, i.e., up to $2^n-1$ constraints. For Pattern~A at $n=30$ this is $\sim 10^9$ and at $n=50$ it is $\sim 10^{15}$, both intractable to enumerate directly.

\paragraph{Sampling strategy.} We solve a \emph{restricted} Core LP over a sampled subset $\mathcal{S}\subseteq\mathcal{F}$: all singletons, all feasible complements $N\setminus\{i\}$, the grand coalition, and $K=10{,}000$ additional random intermediate-size coalitions retained after feasibility filtering. TSP costs $c^{\ast}(S)$ are computed by Held--Karp for $|S|\le 15$ and LKH \citep{helsgaun2000effective} otherwise. The first-stage nucleolus LP
\begin{align}
  \min_{x,\varepsilon}\ \varepsilon\quad\text{s.t.}\quad
  & \sum_i x_i = C(N)_{\text{online}},\qquad
    \sum_{i\in S} x_i - c^{\ast}(S)\le \varepsilon\;\;(S\in\mathcal{S},\ S\ne N) \notag
\end{align}
is solved by CBC via PuLP. If $\varepsilon^{\ast}>0$, $\mathcal{S}$ certifies Core emptiness (the true restricted Core is a subset of the unrestricted Core); the method is one-sided and cannot conclude nonemptiness from a sample.

\paragraph{Scale-up certification (Pattern A seed 123).} After the Phase-4 simulator corrections, the previous 5-case study of Appendix C collapses to the 2 cases at $n=20, 30$ where seed 123 under Pattern A still eludes Theorem~\ref{thm:empty-core}:
\begin{center}
\footnotesize
\begin{tabular}{@{}lccccccc@{}}
\toprule
Pattern & $n$ & Seed & $r$ & $r^{\ast\ast}$ & $\varepsilon^{\ast}$ & Binding sizes & Mechanism \\
\midrule
A & 20 & 123 & 1.077 & 1.112 & $0.77$ & $\{n-2,\,n-1\}$ & near-complement \\
A & 30 & 123 & 1.038 & 1.088 & $0.66$ & $\{n-2,\,n-1\}$ & near-complement \\
\bottomrule
\end{tabular}
\end{center}
The remaining seed-123 scale-up cases at $n=50$ (and at $n\in\{20,30,50\}$ under Pattern~C) now fire Theorem~\ref{thm:empty-core} directly and do not require the restricted-LP certificate.

\paragraph{Main-grid near-complement cases (Remark~\ref{rem:near-complement}).} Table~\ref{tab:near-main} lists the 10 NN main-grid instances classified as near-complement by the first-stage LP's binding-size distribution.

\begin{table}[htbp]
\centering
\footnotesize
\caption{Ten main-grid NN instances in the near-complement class (Remark~\ref{rem:near-complement}). Thm~\ref{thm:empty-core} and Prop~\ref{prop:bondareva-complement} are both vacuous (or the latter inapplicable). Tight binding concentrates at sizes $\{n-2, n-1\}$.}
\label{tab:near-main}
\begin{tabular}{@{}rllrrrrr@{}}
\toprule
$n$ & pattern & seed & $k$ & $r$ & $r^{\ast\ast}$ & $r^{\ast\ast\ast}$ & binding at $n-2, n-1$ \\
\midrule
10 & A        & 123 & 10 & 1.053 & 1.150 & 1.053 & yes \\
10 & B\_heavy & 256 & 10 & 1.126 & 1.138 & 1.061 & yes \\
10 & C        & 123 & 10 & 1.124 & 1.150 & 1.053 & yes \\
15 & A        &  42 & 15 & 1.048 & 1.121 & 1.051 & yes \\
15 & A        &  99 & 15 & 1.051 & 1.151 & 1.043 & yes \\
15 & A        & 123 & 15 & 1.075 & 1.145 & 1.042 & yes \\
15 & B\_heavy &   7 & 15 & 1.185 & 1.194 & 1.046 & yes \\
15 & C        &  42 & 15 & 1.068 & 1.121 & 1.051 & yes \\
15 & C        &  99 & 15 & 1.148 & 1.151 & 1.043 & yes \\
15 & C        & 123 & 15 & 1.141 & 1.145 & 1.042 & yes \\
\bottomrule
\end{tabular}
\end{table}

\paragraph{Main-grid intermediate cases (Observation~\ref{obs:intermediate}).} Table~\ref{tab:inter-binding} gives the full binding-size distribution for the 9 intermediate instances. Binding concentrates on sizes $|S|\in\{3,\ldots,n-3\}$, with no size-$(n-1)$ or size-$(n-2)$ involvement.

\begin{table}[htbp]
\centering
\footnotesize
\caption{Nine main-grid NN instances in the intermediate-coalition class (Observation~\ref{obs:intermediate}). The ``binding sizes'' column lists $|S|=c$ pairs, meaning $c$ constraints of size $|S|$ are tight at the first-stage LP optimum.}
\label{tab:inter-binding}
\begin{tabular}{@{}rllrrrl@{}}
\toprule
$n$ & pattern & seed & $k$ & $r$ & $\varepsilon^{\ast}$ & binding sizes (count) \\
\midrule
 7 & B\_medium & 256 &  5 & 1.436 & 0.028 & $\{1{=}1,\,2{=}1,\,3{=}4,\,4{=}6,\,5{=}2\}$ \\
10 & B\_medium & 123 &  8 & 1.423 & 0.192 & $\{1{=}1,\,2{=}3,\,3{=}6,\,4{=}9,\,5{=}9,\,6{=}5,\,7{=}1\}$ \\
10 & B\_medium & 256 &  8 & 1.358 & 0.217 & $\{1{=}2,\,2{=}2,\,3{=}1,\,4{=}2,\,5{=}1,\,6{=}3,\,7{=}3,\,8{=}1\}$ \\
10 & E        &  42 &  8 & 1.495 & 0.062 & $\{1{=}1,\,2{=}1,\,3{=}5,\,4{=}3,\,5{=}3,\,6{=}2,\,7{=}1\}$ \\
12 & B\_medium &  99 & 10 & 1.395 & 0.744 & $\{3{=}2,\,4{=}2,\,5{=}6,\,6{=}12,\,7{=}12,\,8{=}6,\,9{=}1\}$ \\
12 & B\_medium & 123 & 10 & 1.487 & 0.767 & $\{1{=}1,\,2{=}1,\,3{=}2,\,4{=}2,\,5{=}6,\,6{=}9,\,7{=}7,\,8{=}4,\,9{=}1\}$ \\
12 & B\_medium & 256 & 10 & 1.423 & 0.941 & $\{1{=}1,\,2{=}1,\,6{=}3,\,7{=}9,\,8{=}9,\,9{=}5,\,10{=}1\}$ \\
12 & E        &  42 &  9 & 1.610 & 0.289 & $\{3{=}4,\,4{=}5,\,5{=}5,\,6{=}5,\,7{=}3,\,8{=}1\}$ \\
15 & B\_medium & 256 & 12 & 1.326 & 0.185 & $\{1{=}1,\,2{=}1,\,3{=}3,\,4{=}2,\,7{=}1,\,8{=}5,\,9{=}7,\,10{=}7,\,11{=}2,\,12{=}1\}$ \\
\bottomrule
\end{tabular}
\end{table}

\paragraph{Interpretation.} The near-complement table (Table~\ref{tab:near-main}) confirms that the 10 Remark~\ref{rem:near-complement} cases share the signature $\{n-2, n-1\}$ binding: neither size alone certifies emptiness, but the combined presence of both sizes does. The intermediate table (Table~\ref{tab:inter-binding}) shows that the 9 Observation~\ref{obs:intermediate} cases have binding concentrated on small-to-medium coalition sizes, with $r\in[1.33,1.61]$ and $k<n-1$ in every case. A tight analytic sufficient condition for either class remains open (\S\ref{sec:conclusion}).

\end{document}
